\documentclass[11pt,oneside]{article}
\usepackage[a4paper,margin=1in]{geometry}
\usepackage[T1]{fontenc}
\usepackage[utf8]{inputenc}
\usepackage{lmodern}
\usepackage{microtype}
\usepackage{amsmath,amssymb,mathtools}
\usepackage{natbib}
\usepackage{booktabs,longtable,array,multirow}
\usepackage{enumitem}
\usepackage{xcolor}
\usepackage{hyperref}
\usepackage{graphicx}
\newcolumntype{P}[1]{>{\raggedright\arraybackslash}p{#1}}
\hypersetup{colorlinks=true,linkcolor=blue!50!black,citecolor=blue!50!black,urlcolor=blue!50!black}
\setlist{itemsep=2pt,topsep=4pt}
\newcommand{\E}{\mathbb{E}}
\newcommand{\Prb}{\mathbb{P}}
\newcommand{\R}{\mathbb{R}}
\newcommand{\logit}{\operatorname{logit}}
\newcommand{\sigm}{\operatorname{sigm}}
\newcommand{\argmin}{\operatorname*{arg\,min}}
\newcommand{\argmax}{\operatorname*{arg\,max}}
\newcommand{\given}{\,|\,}
\title{Forecasting bank product opportunities and allocating scarce sales capacity\\\large A literature survey and product proposal}
\author{Research design note}
\date{8 October 2026}

\begin{document}
\maketitle
\begin{abstract}
This document specifies a decision system for a commercial or corporate bank that has roughly one million business clients, only 36 monthly observations per client, internal transaction records, client attributes, and dated sales contacts and outcomes. The system must identify which client--product pairs are likely to need attention soon and decide which contacts are worth scarce salesperson time. The two decisions require different statistical quantities. A forecast estimates a future bank-observed event. An outreach model estimates the incremental outcome caused by a contact. An allocation model converts those estimates into a capacity-constrained portfolio. Treating these quantities as one score creates a predictable error: the clients most likely to transact are not necessarily the clients for whom a conversation changes the bank's result.

The survey examines ten primary papers in a common notation. It derives pooled and global forecasting, discrete-time survival analysis, the beta-geometric\/negative-binomial model, a sequential cross-selling latent-maturity model, gradient boosting, DeepAR probabilistic forecasting, positive--unlabeled learning, uplift modelling, and observational policy learning. Each method is linked to the bank setting and its assumptions are tested against the short panel, selective observation across banks, repeated products, and historical sales selection. The proposal is a product-specific, pooled discrete-time hazard model with hierarchical effects and a calibrated count component, a separately estimated outreach-effect model, and a constrained policy optimizer. It uses global models to borrow strength across clients, survival likelihoods for timing, product-history features for cross-selling, and validated causal policy recommendations after an explicit contact experiment. The proposed evaluation uses chronological backtesting, constructed heuristic baselines, clustered uncertainty, and a prospective randomized pilot. No bank-specific performance claim is made without client data.
\end{abstract}

\clearpage
\begingroup
\footnotesize
\setcounter{tocdepth}{2}
\tableofcontents
\endgroup
\clearpage

\section{The business decision}
\subsection{A precise statement of the problem}

Let $i\in\{1,\ldots,N\}$ index business clients, $j\in\{1,\ldots,J\}$ index bank products, and $t\in\{1,\ldots,36\}$ index calendar months. A product is defined at the level at which a commercial decision is made: for example, an FX hedge, a revolving credit facility, a term loan, a deposit or liquidity service, a payment or collections service, or a trade-finance service. The definition should be agreed with product owners before modelling. A vague target such as ``customer activity'' joins events with different commercial consequences and makes a score difficult to act on.

Write $Y_{ijt}^{(H)}=1$ when client $i$ has a qualifying bank-observed event for product $j$ during the $H$ months after the information cutoff at month $t$. A qualifying event might be a new product opening, a funded draw, an FX transaction above a product-specific threshold, a material increase in a facility, or another event for which a salesperson can plausibly help. Let $C_{it}$ be the client information available by the cutoff: lagged transaction summaries, current products, account tenure, industry, geography, legal-entity relations, risk and eligibility fields, and any external feature whose publication date is no later than the cutoff. The forecasting target is
\begin{equation}
 p_{ijt}^{(H)}=\Prb\left(Y_{ijt}^{(H)}=1\given C_{it}\right).
 \label{eq:forecast-target}
\end{equation}
 This is a conditional probability of an event visible to the bank. It is not the probability that the client has a need somewhere in the financial system.

The product may also report a count. Let $N_{ijt}^{(H)}$ be the number of qualifying transactions in the horizon. A forecast of its conditional mean, $m_{ijt}^{(H)}=\E[N_{ijt}^{(H)}\given C_{it}]$, describes transaction intensity. Separate variables are needed for dollar amounts, balances, fees and notional volume; their units and economic meaning differ. A binary opportunity flag is often easier to calibrate; a count or value distribution is useful for ranking large opportunities. The document uses both, with the binary hazard as the primary short-panel target.

The historical contact indicator $W_{it}$ records assignment to a defined sales protocol after the information cutoff. Assignment begins the outcome window; delivery occurs within its specified initial response period. Records must distinguish assignment, attempted delivery and completed contact. Define potential outcomes $Y_{it}^{(H)}(1)$ and $Y_{it}^{(H)}(0)$ under contact and no contact. For a contribution margin $M_{ijt}^{(H)}(w)$, the incremental value of a contact can be written
\begin{equation}
 \tau_{ijt}=\E\left[M_{ijt}^{(H)}(1)-M_{ijt}^{(H)}(0)\given C_{it}\right],
 \label{eq:tau}
\end{equation}
 where margin excludes the separately recorded outreach cost $k_{it}$. Net action value is $v_{ijt}=\tau_{ijt}-k_{it}$. Equation \eqref{eq:tau} is a causal estimand. A high value of $p_{ijt}^{(H)}$ does not imply a high value of $\tau_{ijt}$: a client may transact without a conversation, or a conversation may redirect business from another bank.

Finally, let $a_{it}\in\{0,1,\ldots,J\}$ denote the action assigned to a client, and let $h_{it}$ be the required salesperson hours. At a fixed cutoff $t$, with monthly capacity $B$, a simple policy solves
\begin{equation}
 \max_{a}\sum_i v_{it}(a_{it})
 \quad\text{subject to}\quad
 \sum_i h_{it}(a_{it})\le B,
 \quad a_{it}\in\mathcal A_{it},
 \label{eq:allocation}
\end{equation}
 where $v_{it}(a)$ is expected incremental contribution and $\mathcal A_{it}$ contains eligibility, ownership, conflict, and contact-frequency constraints. If calls compete only for a common pool of hours, \eqref{eq:allocation} is a knapsack problem. If a salesperson can contact several clients in a territory, or a client must receive a coordinated product bundle, the constraints become a mixed-integer program. The statistical model supplies values and uncertainty; it does not remove the business constraints.

\subsection{The unit of prediction}

The production table should have one row per client, product, cutoff month, and horizon. It must include the feature timestamp and the label end date. A client can have several rows at the same cutoff because FX, lending, deposits, and payments are distinct opportunities. A product event should be marked once in the horizon even if it contains many transactions, unless the commercial action is explicitly volume expansion. Product rows are therefore not independent observations merely because they are numerous. They share a client, a relationship manager, a month, and macroeconomic conditions.

The target is a forward-looking event. For a cutoff at the end of March and a three-month horizon, features must be frozen at March 31 and the label covers April through June. A contact recorded in April cannot be a March feature. A web page first published in May cannot be used as March information merely because it is downloaded today. This point is operational rather than cosmetic: leakage turns a historical reconstruction into an impossible prediction task.

\subsection{A concrete illustration}

Consider a hypothetical product for which the event probability under usual service is $0.80$ for client A and $0.20$ for client B. Suppose the additional call leaves A's probability at $0.80$ but raises B's to $0.50$. With equal event margins and contact costs, A has no incremental response and B has a $0.30$ probability increase. Ranking by usual-service event probability puts A first; ranking by incremental response puts B first. These are assumed numbers for illustration, not effects inferred from the bank's records. The product must distinguish the two quantities.

The bank should therefore expose three fields: the predicted event probability, the estimated incremental value of each eligible action, and the constraints or reason codes that produced the final assignment. Forecast reason codes are descriptive associations such as recent payment growth or an approaching facility maturity. Causal reason codes require a validated treatment model and should be labelled as estimated treatment heterogeneity, not as a mechanistic explanation.


\section{Data, estimands, and identification}
\subsection{Information available to the bank}

For each cutoff month, construct features in five groups. Transaction features include recency, frequency, amount, volatility, trend, seasonality, payment corridors, currencies, counterparties where permitted, and product-specific utilization. Relationship features include product holdings, tenure, limits, balances, prior service quality, ownership, documented pipeline, and contact history. Client features include industry, size bands, legal structure, geography, address stability, and risk or eligibility variables. Calendar features include month, holidays, rates, commodity or sector indicators, and bank-wide product changes. External web features may describe a registered activity or announced facility only when the source and publication timestamp are stored; the text itself is not evidence that the client will transact.

The feature vector must include missingness indicators. A missing balance can mean that an account does not exist, that a feed arrived late, or that the client opted out. Those states have different meanings. Imputing them to a common zero silently turns data availability into a behavioural assertion.

\subsection{Short panels and dependence}

Thirty-six months gives at most 36 cutoff rows per client and fewer after lagging, maturity, and horizon censoring. One million clients produce many cross-sectional examples, so a model can estimate shared effects and rare categorical patterns. They do not supply one million independent realizations of a rate shock or a regulatory event. The effective number of calendar regimes is near the number of distinct months, not the number of rows. Client rows are also correlated within legal-entity groups, relationship managers, regions, products, and months. Confidence intervals and validation splits must reflect these clusters.

The short panel favours low-dimensional lag summaries, hierarchical shrinkage, and direct one- or three-month forecasts. A separate model for each client's time series would estimate too many parameters from too little history. A very deep model can memorize client identifiers or month effects while appearing accurate in a random row split. The baseline should pool clients; local models and high-capacity challengers earn their place through chronological evaluation.

\subsection{What the bank can identify}

Let $D_{it}$ indicate a total-market product need and $S_{it}$ indicate that the bank captures it. The bank records $Y_{it}=D_{it}S_{it}$ for its own transactions. Even with perfect measurement of $q(x)=\Prb(Y=1\given X=x)$, the two factors are not identified. For example, $q(x)=0.10$ is compatible with $(\Prb(D=1\given x),\Prb(S=1\given D=1,x))=(0.20,0.50)$ and $(0.50,0.20)$. The observed product is the same. No classifier trained only on internal transactions can choose between these worlds.

If an independent source supplies a defensible bound $c_L\le\Prb(S=1\given D=1,x)\le c_U$, with $0<c_L\le c_U\le1$, then, when the bounds are compatible with $q(x)$,
\begin{equation}
 \frac{q(x)}{c_U}\le \Prb(D=1\given x)\le \min\!\left(1,\frac{q(x)}{c_L}\right).
 \label{eq:wallet-bounds}
\end{equation}
 Equation \eqref{eq:wallet-bounds} is a sensitivity interval, not an estimate of wallet share. A web signal can help predict a bank-visible opportunity or provide a prior for $c(x)$, but it cannot by itself identify the unseen denominator. A survey, consented data-sharing arrangement or representative industry panel that measures otherwise hidden need and capture can supply additional information. A randomized offer identifies effects on the observed outcome; randomization alone does not reveal the unobserved total-market denominator.

\subsection{Missing labels and selective sales records}

For an eligible client--product pair with complete observation, no qualifying transaction is a valid negative for the own-bank target, even if the client transacts elsewhere. It is not a negative label for total-market need. A feed gap or an unfinished outcome horizon is an unknown label, while ineligibility removes the pair from the relevant risk set. These distinctions should be represented explicitly. Sales outcomes are also selected because bankers choose whom to contact. A contacted client can have high natural propensity even if the contact has little effect, so a naive conversion comparison confounds selection with treatment. The forecasting target can use all completely observed eligible rows. The causal target additionally needs assignment, eligibility, timing and a defined outcome comparison.

The first product release should use a forecast target that the bank measures reliably and present total-market need as a scenario or sensitivity range. It should not advertise ``wallet share'' from internal data alone.


\section{What the literature contributes}
The ten papers reviewed here answer complementary questions. Montero-Manso and Hyndman explain why many short histories can support one forecasting function. Singer and Willett, and their later multiple-spell paper, explain how to model the time until an event. Fader, Hardie, and Lee supply a parsimonious model of repeat purchases and permanent inactivity. Li, Sun, and Wilcox model dependencies between banking products. Friedman and Salinas and colleagues provide nonlinear predictive methods. Gutierrez and Gérardy distinguish response prediction from treatment response. Athey and Wager connect treatment effects to policy choice. Elkan and Noto expose the assumptions needed to learn from selectively observed positives. The survey is a focused methodological synthesis; it is not an exhaustive catalogue of commercial lead-scoring systems.

For each paper, the discussion reconstructs the mathematical ingredients used in the proposal, including the likelihood or loss, fitting procedure, prediction rule, and the assumptions needed to interpret it. Additional bank-specific derivations are identified as adaptations. Published empirical results establish what the authors studied, not performance at this bank. No numerical result from a consumer, retail-demand, or simulation example is treated as a corporate-banking performance estimate.

\subsection{Global forecasting: pooling does not require identical clients}

\citet{montero2021global}, Sections 2--3, distinguish a local learning algorithm $A_L$, applied separately to each series, from a global algorithm $A_G$, fitted across the collection. If $X_i=(y_{i1},\ldots,y_{iT})$ denotes the available history, a local forecast is $f_i(X_i)$, where $f_i=A_L(X_i)$; a global forecast is $g(X_i)$, where $g=A_G(X_1,\ldots,X_N)$. Shared parameters do not require equal forecasts. Different histories, industries and product holdings can produce different predictions from one function.

Their Proposition 1 concerns representation on a finite set of histories. To reconstruct its main argument, first suppose a global function $g$ has been supplied. A local algorithm that returns this same function for every input reproduces its forecasts. Conversely, apply a deterministic local algorithm to each observed history and form the finite set of pairs $(X_i,f_i(X_i))$. If two complete histories are identical, the deterministic algorithm returns the same fitted function and the same prediction for both. Thus these pairs define a function on the distinct observed histories, which can be extended outside that finite set in any manner. A sufficiently flexible global function can reproduce those particular local predictions. This construction proves representability, not that fitting the global function will discover the right extension or generalize well.

The use of the complete history matters. If a global model receives only the last three lags, two clients with the same last three observations but different older histories can have different local predictions. A function of those three lags alone cannot reproduce both. The paper discusses this restriction in Section 2.2. Adding longer history or appropriate characteristics can distinguish the clients; adding an identifier can also distinguish them, but may encourage memorization. The theorem itself does not require client identifiers and does not justify unrestricted identifier embeddings in a short panel.

\subsubsection{Why sharing parameters can reduce estimation error}

A finite-class calculation makes the statistical argument precise. Consider $K$ independent series, each supplying $n$ independent loss observations in $[0,1]$. This deliberately stronger toy assumption makes every step transparent. For a fixed predictor $f$, write $Z_r=\ell_r(f)-\E\ell_r(f)$ for $r=1,\ldots,Kn$. The bounded-variable moment inequality gives $\E e^{sZ_r}\le e^{s^2/8}$. Markov's inequality and independence then imply
\[
 \Prb\left(\sum_r Z_r>Kn\epsilon\right)
 \le \exp(-sKn\epsilon)\prod_r\E e^{sZ_r}
 \le \exp(-sKn\epsilon+Kn s^2/8).
\]
Minimizing the exponent at $s=4\epsilon$, and applying the same argument to $-Z_r$, gives
\[
 \Prb\bigl(|\widehat R(f)-R(f)|>\epsilon\bigr)
 \le 2e^{-2Kn\epsilon^2}.
\]
For a finite collection $\mathcal F$, the probability that any predictor exceeds this error is no greater than the sum of those probabilities. Setting this sum to $\delta$ gives
\begin{equation}
 \sup_{f\in\mathcal F}|\widehat R(f)-R(f)|
 \le \sqrt{\frac{\log|\mathcal F|+\log(2/\delta)}{2Kn}}
 \quad\text{with probability at least }1-\delta.
 \label{eq:global-bound}
\end{equation}
A local system can choose a separate predictor from each class $\mathcal H_i$, so its collection has size $\prod_i|\mathcal H_i|$ and logarithmic complexity $\sum_i\log|\mathcal H_i|$. A global system chooses one predictor from $\mathcal G$ and incurs $\log|\mathcal G|$. For equally sized local classes, the former complexity grows as $K\log|\mathcal H|$. Pooling can therefore permit a richer shared function without paying separately for its complexity on every series. This reconstructs the finite-class comparison in the paper's Sections 3.1--3.3; the paper discusses an effective sample size for within-series dependence.

The calculation does not supply a numerical confidence bound for the bank. Losses share clients, business groups and calendar shocks; margin loss is not bounded without further assumptions; and modern model classes are not finite enumerations. In particular, $36$ million rows cannot be substituted for $Kn$ as though they were independent. The defensible implication is to compare shared and partially pooled models, while estimating uncertainty using the actual dependence structure.

For this bank, a common coefficient can learn from many importers that a change in foreign-currency payments precedes an FX transaction. Industry interactions allow exporters to behave differently. A hierarchical coefficient allows a small industry to borrow information from larger related industries. The paper's forecasting benchmarks do not establish the size of this benefit for bank sales; that remains an out-of-time comparison.

\subsection{BG/NBD: repeat activity and latent permanent inactivity}

\citet{fader2005counting}, Sections 2--3 and Appendix A, model repeat transactions after an initial acquisition. While active, an individual transacts according to a Poisson process with rate $\lambda>0$. Immediately after each repeat transaction the individual becomes permanently inactive with probability $p\in(0,1)$. Across individuals,
\[
 \lambda\sim\operatorname{Gamma}(r,\text{rate }\alpha),\qquad
 p\sim\operatorname{Beta}(a,b),\qquad \lambda\perp p,
\]
where all four population parameters are positive. The gamma density is $\alpha^r\lambda^{r-1}e^{-\alpha\lambda}/\Gamma(r)$, and the beta density is $p^{a-1}(1-p)^{b-1}/B(a,b)$, with $B(a,b)=\Gamma(a)\Gamma(b)/\Gamma(a+b)$.

For an observation window of length $T$, let $x$ be the number of repeat transactions and $t_x$ the time of the last repeat, measured from acquisition. The original process starts active after acquisition. In particular, with $x=0$ it has not yet encountered a repeat-transaction dropout opportunity. There are two mutually exclusive explanations for a history with $x>0$: the individual remains active through $T$, or drops out immediately after transaction $x$. Up to parameter-independent terms, the conditional likelihood is
\begin{equation}
 L(\lambda,p\mid x,t_x,T)
 = (1-p)^x\lambda^x e^{-\lambda T}
 +\mathbf 1_{\{x>0\}}p(1-p)^{x-1}\lambda^x e^{-\lambda t_x}.
 \label{eq:bgnbd-conditional}
\end{equation}
The first term requires surviving all $x$ dropout opportunities and the observed arrival history through $T$. The second requires surviving the first $x-1$ opportunities and dropping out at the last arrival. There is no further no-arrival requirement after permanent dropout.

\subsubsection{Integrating the heterogeneity}

For any $u\ge0$, direct gamma integration gives
\[
 \int_0^\infty\lambda^xe^{-u\lambda}
 \frac{\alpha^r\lambda^{r-1}e^{-\alpha\lambda}}{\Gamma(r)}\,d\lambda
 =\frac{\alpha^r\Gamma(r+x)}{\Gamma(r)(\alpha+u)^{r+x}}.
\]
Similarly, integrating $(1-p)^x$ against the beta density gives $B(a,b+x)/B(a,b)$, while integrating $p(1-p)^{x-1}$ gives $B(a+1,b+x-1)/B(a,b)$. Hence the marginal likelihood is $A+D$, where
\begin{align}
 A&=\frac{B(a,b+x)}{B(a,b)}
       \frac{\alpha^r\Gamma(r+x)}{\Gamma(r)(\alpha+T)^{r+x}},\label{eq:bgnbd-A}\\
 D&=\mathbf 1_{\{x>0\}}\frac{B(a+1,b+x-1)}{B(a,b)}
       \frac{\alpha^r\Gamma(r+x)}{\Gamma(r)(\alpha+t_x)^{r+x}}.\label{eq:bgnbd-D}
\end{align}
Estimate the population parameters by maximizing $\sum_i\log(A_i+D_i)$, using logarithmic parameterizations for positivity and a log-sum-exp calculation for numerical stability. Clients contribute different sufficient summaries $(x_i,t_{x_i},T_i)$; the four population parameters are shared.

Bayes' rule gives
\begin{equation}
 P_{\rm active}=\frac{A}{A+D}
 =\left[1+\mathbf1_{\{x>0\}}\frac{a}{b+x-1}
       \left(\frac{\alpha+T}{\alpha+t_x}\right)^{r+x}\right]^{-1},
 \label{eq:bgnbd-alive}
\end{equation}
where the ratio expression is used only for $x>0$; when $x=0$, $D=0$. Long time since the last repeat raises $D/A$ and lowers the inferred probability of remaining active. More recent activity raises it. The $x=0$ result exposes a substantive assumption: the original model cannot infer pre-first-repeat dropout. This can be implausible for a newly onboarded but unused bank account.

\subsubsection{A self-contained future-count forecast}

Conditional on being active at $T$, conjugacy gives independent posteriors
\[
 \lambda\mid\text{active, history}\sim\operatorname{Gamma}(r+x,\alpha+T),
 \quad p\mid\text{active, history}\sim\operatorname{Beta}(a,b+x).
\]
For fixed $(\lambda,p)$, dropout events form a thinning of the active Poisson process with rate $\lambda p$. The probability of still being active at future elapsed time $u$ is $e^{-\lambda pu}$. Expected transaction intensity is therefore $\lambda e^{-\lambda pu}$, and integrating over $0\le u\le H$ gives $(1-e^{-\lambda pH})/p$. Averaging first over the gamma posterior, using its Laplace transform, gives the exact finite-horizon forecast
\begin{equation}
 \E[N(T,T+H)\mid\text{history}]
 =\frac{P_{\rm active}}{B(a,b+x)}
 \int_0^1 p^{a-2}(1-p)^{b+x-1}
 \left[1-\left(\frac{\alpha+T}{\alpha+T+pH}\right)^{r+x}\right]dp.
 \label{eq:bgnbd-forecast}
\end{equation}
This one-dimensional integral is directly computable. Although its first factor contains $p^{a-2}$, the bracket is asymptotic to $(r+x)pH/(\alpha+T)$ near zero, so the complete integrand is integrable for every $a>0$. It avoids requiring a reader to evaluate an unexplained special function.

For completeness, the familiar closed form in the paper follows when $a>1$. Put $q=r+x$, $c=\alpha+T$ and $z=H/(c+H)$. Split the integral into its two terms and use
\[
 {}_2F_1(A,B;C;z)=\frac{1}{B(B,C-B)}
 \int_0^1 u^{B-1}(1-u)^{C-B-1}(1-zu)^{-A}du,
\]
which can be verified by expanding $(1-zu)^{-A}=\sum_{n\ge0}(A)_nz^nu^n/n!$ and integrating each beta moment; $(A)_n=A(A+1)\cdots(A+n-1)$. The substitution $u=1-p$ then yields
\begin{align}
 \E[N(T,T+H)\mid\text{history}]
 &=P_{\rm active}\frac{a+b+x-1}{a-1}\nonumber\\[-2pt]
 &\quad\times\left[1-\left(\frac{c}{c+H}\right)^q
 {}_2F_1(q,b+x;a+b+x-1;z)\right].
 \label{eq:bgnbd-hypergeom}
\end{align}
The integral form remains the unambiguous definition at parameter values where this split-integral derivation is unavailable.

\subsubsection{What transfers to monthly bank data}

The useful insight is a separation between frequency while active and uncertainty about continued activity. Recency and frequency can be informative with very short histories. However, permanent inactivity with this bank is not permanent absence of financial need: the client may move activity to another bank and later return. Product adoption, seasonal corporate flows and repeated facility draws also differ from the original repeat-purchase setting. The paper's customer-base illustration cannot validate these assumptions for business clients.

There is an additional measurement constraint. An exact last transaction time $t_x$ is not observed in a monthly aggregate. Inserting month-end as its exact time changes the likelihood. With exact dated transactions recovered from the underlying system, the original BG/NBD can be a benchmark. With only monthly counts, use a genuinely discrete model or integrate the unobserved timestamps under an explicitly specified observation model. The next section gives a discrete analogue. It is a new adaptation, not the original BG/NBD under a new name.


\section{Forecasting event timing and transaction counts}
\subsection{Single-event discrete-time survival analysis}

\citet{singer1993time}, especially the development of the hazard and person-period likelihood and their Equation (3), explain how logistic regression can model the timing of an event observed in discrete intervals. Fix a client $i$ and product $j$. Let $T_{ij}$ be the month of the first qualifying event after a clearly defined origin, such as becoming eligible for a product. At the beginning of month $t$, the pair is at risk if it is eligible, observable and has not yet experienced that event. The hazard is
\begin{equation}
 h_{ijt}=\Prb(T_{ij}=t\given T_{ij}\ge t,\mathcal H_{ij,t-1}),
 \label{eq:hazard}
\end{equation}
where $\mathcal H_{ij,t-1}$ contains information known before that month. Conditional on a specified covariate path, the chain rule gives
\begin{equation}
 S_{ij}(t)=\Prb(T_{ij}>t)=\prod_{u=1}^{t}(1-h_{iju}),\qquad
 \Prb(T_{ij}=t)=h_{ijt}\prod_{u<t}(1-h_{iju}).
 \label{eq:survival-product}
\end{equation}
The product is a product of successive conditional probabilities. It does not require unconditional independence of a client's event indicators. If future covariates are random, the unconditional survival probability averages this conditional product over their possible paths; it is generally not the product of separately averaged hazards.

Let $c_{ij}$ be the last observable interval and $\delta_{ij}=1$ if the event was observed by that interval. An event at $c_{ij}$ contributes $h_{ijc_{ij}}\prod_{u<c_{ij}}(1-h_{iju})$; a right-censored observation contributes $\prod_{u\le c_{ij}}(1-h_{iju})$. Writing $r_{ijt}$ for membership in the observed risk set, and $y_{ijt}$ for an event in that row, combines both cases:
\begin{equation}
 L=\prod_{i,j,t}\left\{h_{ijt}^{y_{ijt}}(1-h_{ijt})^{1-y_{ijt}}\right\}^{r_{ijt}}.
 \label{eq:first-event}
\end{equation}
Rows after first adoption are absent from a first-adoption model. A currently owned product can still be eligible for a separately defined expansion event. Treating its ownership as an additional first adoption would give the model the wrong target.

Ordinary right-censoring calculations require observation to end independently of the event process, conditional on the modeled information. A feed outage must not be recorded as a zero. A relationship closure may be an economically informative competing outcome, rather than independent censoring. For a target defined as an own-bank transaction over a fully observable horizon, a closed relationship with verified no transaction can be a genuine zero; for a target requiring continued eligibility, closure needs its own transition or a clearly changed target. These choices belong in the label specification.

\subsubsection{Logistic estimation and its derivatives}

Use $\sigm(z)=1/(1+e^{-z})$ and set
\begin{equation}
 h_{ijt}=\sigm(\eta_{ijt}),\quad
 \eta_{ijt}=\alpha_j(k_{ijt})+x_{ij,t-1}^{\mathsf T}\beta_j
             +b_{j,g(i)}+u_{ij}+s_j(t).
 \label{eq:hazard-model}
\end{equation}
Here $k_{ijt}$ is elapsed duration, $\alpha_j(k)$ is the baseline duration effect, $g(i)$ is an industry or other segment, and $s_j(t)$ contains calendar features available for forecasting. Product index $j$, duration $k$, and calendar month $t$ play different roles. A separate unconstrained coefficient for each historical calendar month cannot forecast a new month. Month-of-year effects, known holidays and documented macroeconomic scenarios can. With three years, annual seasonality remains weakly established and should be shrunk strongly or omitted if validation does not support it.

For a design row $z$ and coefficient vector $\theta$, the negative log-likelihood is
\[
 \mathcal L(\theta)=\sum_r\{\log(1+e^{z_r^{\mathsf T}\theta})-y_rz_r^{\mathsf T}\theta\}
                       +\tfrac12\theta^{\mathsf T}\Lambda\theta,
\]
where the sum includes only risk-set rows and $\Lambda$ encodes a declared quadratic penalty. Since $\partial h/\partial\eta=h(1-h)$,
\begin{align}
 \nabla\mathcal L&=Z^{\mathsf T}(h-y)+\Lambda\theta,\label{eq:row-loglik}\\
 \nabla^2\mathcal L&=Z^{\mathsf T}\operatorname{diag}\{h_r(1-h_r)\}Z+\Lambda.
\end{align}
A Newton step solves this Hessian system for the coefficient update. Equivalently, iteratively reweighted least squares uses weights $h_r(1-h_r)$ and working responses $z_r^{\mathsf T}\theta+(y_r-h_r)/[h_r(1-h_r)]$. Stable logistic and linear-algebra routines are necessary for rare events; literal exponentiation of large predictors and division by numerically zero weights are not acceptable implementations. Penalty strengths and any numerical modification require validation and a recorded rationale.

A complementary log-log link is a useful local extension, not needed for the logistic model. Suppose the event intensity is constant at $\lambda$ during an interval of length $\Delta$. Solving $S'(u)=-\lambda S(u)$ with $S(0)=1$ gives $S(\Delta)=e^{-\lambda\Delta}$. Thus $h=1-e^{-\lambda\Delta}$ and $\log[-\log(1-h)]=\log\lambda+\log\Delta$. This provides an exposure offset when intervals differ in length. For the bank's equally spaced monthly data, a logistic hazard is the simpler initial choice.

\subsection{Repeated events: the multiple-spell construction}

\citet{willett1995multiple}, particularly their likelihood development through Equation (19) and the person-spell-period construction, extend the single-event approach to repeated spells. Their application concerns alternating spells in and out of teaching. The transferable idea is to define which state a subject occupies, which transition ends a spell, and which history is available at the next spell's origin. Their application is not a bank-sales validation.

For FX transactions, define a spell as the wait from one qualifying event month to the next. For a client with monthly event indicators $(0,1,0,0,1)$, the first observed spell has duration two and the second has duration three, if the first month is a known risk origin. If monitoring starts mid-spell, its true duration is unknown; either condition on delayed entry with known duration or distinguish this left-truncated history from a fully observed spell. A transaction in the same month as another transaction does not create an observed one-day spell: monthly indicators cannot resolve it. Counts require the count model below.

Let $s$ index spells and $k$ their elapsed months. With $\mathcal H_{ijs,k-1}$ containing previous spell lengths, previous outcomes and available time-varying features, define $h_{ijsk}=\Prb(\text{spell }s\text{ ends at }k\mid\text{not yet ended},\mathcal H_{ijs,k-1})$. The conditional likelihood is
\begin{equation}
 L_i(u_i)=\prod_{j,s,k}\left[h_{ijsk}(u_i)^{y_{ijsk}}
                          (1-h_{ijsk}(u_i))^{1-y_{ijsk}}\right]^{r_{ijsk}}.
 \label{eq:recurrent}
\end{equation}
Spell indicators and interactions allow first and subsequent events to have different baselines. A shared unobserved client effect introduces dependence across spells; its marginal likelihood is $\int L_i(u)f(u)du$. Without such a model, clustered uncertainty can accommodate some within-client dependence in estimation, but does not repair an incorrect conditional mean. The Bernoulli-shaped likelihood is not permission to count every person-month as an independent client in a standard-error calculation.

Distinct products need not be mutually exclusive competing risks. A client can adopt a liquidity service and execute FX in the same month. Product-specific hazards permit both; their joint distribution needs a dependence model only when a joint probability or portfolio uncertainty is required. Competing-risk probabilities are appropriate only for genuinely mutually exclusive transitions under the chosen event definition.

\subsubsection{From monthly hazards to a sales horizon}

For a cutoff $t$, define $h_{ij,t+k\mid t}$ as the conditional probability of the first future event in month $t+k$, given no earlier future event and the information at $t$. Then
\begin{equation}
 p_{ijt}^{(H)}=1-\prod_{k=1}^{H}(1-h_{ij,t+k\mid t}).
 \label{eq:horizon-hazard}
\end{equation}
For example, hazards $(0.10,0.20,0.10)$ imply $1-0.9\times0.8\times0.9=0.352$. Multiplying survival probabilities answers ``at least one event,'' not expected event count. Forecasts of future covariates must be integrated or generated prospectively; the realized April balance cannot be used in a forecast issued at the end of March.

A practical alternative is a direct model for each horizon, $\Prb(Y_{ijt}^{(H)}=1\mid C_{it})=\sigm(f_{jH}(C_{it}))$. It avoids recursive covariate forecasting but does not automatically guarantee $p^{(1)}\le p^{(3)}$. A common sequence of nonnegative conditional hazards ensures this nesting. Compare the alternatives using calibrated horizon probabilities rather than assuming one is superior.

\subsection{Partial pooling and sparse clients}

A million fixed client intercepts can absorb noise from short histories. A normal prior $u_{ij}\sim N(0,\sigma_{u,j}^2)$ penalizes implausibly large client effects and supplies a population prediction for new clients. The direction of shrinkage follows exactly in a Gaussian approximation: if $z_i\mid u_i\sim N(u_i,s_i^2)$, then completing the square in
\[
 -\frac{(z_i-u_i)^2}{2s_i^2}-\frac{u_i^2}{2\sigma_u^2}
 =-\frac12(s_i^{-2}+\sigma_u^{-2})
 \left[u_i-\frac{\sigma_u^2}{\sigma_u^2+s_i^2}z_i\right]^2+\text{constant}
\]
gives posterior mean $\sigma_u^2z_i/(\sigma_u^2+s_i^2)$ and variance $(s_i^{-2}+\sigma_u^{-2})^{-1}$. Sparse clients have large $s_i^2$ and shrink more. This closed form illustrates the mechanism; logistic random-effect posteriors are not exactly Gaussian and require integration or an explicitly checked approximation.

The initial bank model should pool industry and product effects and use observed recency/frequency features. Client random effects are an optional extension whose computational cost and cold-start behavior must be justified. Predictions for a new industry or client should integrate an appropriate group-effect distribution, rather than assign an arbitrary trained identifier. Regularization strengths are selected on past-to-future validation, not inherited as universal constants.

\subsection{Counts, positive amounts, and a discrete inactivity model}

A Poisson count with mean $\lambda$ has probability $e^{-\lambda}\lambda^n/n!$ and variance equal to its mean. Corporate transaction counts often require a more variable distribution. Mixing this Poisson likelihood with $\lambda\sim\operatorname{Gamma}(k,\text{rate }k/\mu)$ gives
\begin{align}
 \Prb(N=n)
 &=\int_0^\infty\frac{e^{-\lambda}\lambda^n}{n!}
   \frac{(k/\mu)^k\lambda^{k-1}e^{-k\lambda/\mu}}{\Gamma(k)}d\lambda\nonumber\\
 &=\frac{\Gamma(n+k)}{\Gamma(k)n!}
       \left(\frac{k}{k+\mu}\right)^k
       \left(\frac{\mu}{k+\mu}\right)^n.
 \label{eq:nb}
\end{align}
By iterated expectation, $\E N=\E\lambda=\mu$; by total variance, $\operatorname{Var}(N)=\E\lambda+\operatorname{Var}(\lambda)=\mu+\mu^2/k$. A log link for $\mu$ keeps the mean positive. As $k\to\infty$, heterogeneity vanishes and the Poisson limit returns.

For a hurdle model, let $p=\Prb(N>0)$ and use a zero-truncated negative binomial for positive counts:
\[
 \Prb(N=0)=1-p,\quad
 \Prb(N=n>0)=p\frac{f_{\rm NB}(n;\mu,k)}{1-f_{\rm NB}(0;\mu,k)}.
\]
Consequently $\E N=p\mu/[1-f_{\rm NB}(0)]$. It would be wrong to fit an untruncated count likelihood only to positive counts and retain its original mean formula. For positive monetary amounts $V$, a separate lognormal model can use $\log V\sim N(\mu_V,\sigma_V^2)$; completing the square in the normal moment integral gives $\E[V\mid V>0]=e^{\mu_V+\sigma_V^2/2}$. Counts, dollar amounts, balances and notionals must remain separate quantities. FX notional is not bank revenue.

A discrete analogue of the BG/NBD idea can be written without inventing event dates. Assume an active client generates a monthly Poisson count with rate $\lambda$, and, after a positive month, becomes inactive with probability $q$. Inactive clients generate zero forever. This is an explicitly different month-end dropout mechanism. Conditional on $(\lambda,q)$, let $a_m$ and $d_m$ be joint probabilities of the observed history through month $m$ and active or inactive status. Starting with $a_0=1,d_0=0$, and writing $f_m=e^{-\lambda}\lambda^{n_m}/n_m!$,
\begin{align}
 a_m&=a_{m-1}f_m\left(1-q\mathbf1_{\{n_m>0\}}\right),\\
 d_m&=d_{m-1}\mathbf1_{\{n_m=0\}}
        +a_{m-1}f_m q\mathbf1_{\{n_m>0\}}.
 \label{eq:monthly-forward}
\end{align}
The likelihood is $a_T+d_T$. Integrating it against gamma and beta mixing densities gives a population model, evaluable with bounded quadrature or another checked integration method. Each recursion is just summation over the two possible previous states. Probability of current activity is the integrated $a_T$ divided by the integrated likelihood. Permanent inactivity is still a strong assumption; return to activity would require another transition parameter. For the initial release, recency features and the hazard/hurdle model are simpler than fitting this latent-state extension.

\subsection{Features that preserve the information cutoff}

At cutoff $t$, compute summaries only from recorded, available months. For a complete lookback of $L$ months, examples are
\begin{align*}
 \operatorname{recency}_{ijt}&=t-\max\{u\le t:N_{iju}>0\},\\
 \operatorname{frequency}_{ijt}(L)&=\sum_{u=t-L+1}^t\mathbf1_{\{N_{iju}>0\}},\\
 \operatorname{growth}_{ijt}(2m)&=
 \log\left(1+\frac1m\sum_{u=t-m+1}^{t}N_{iju}\right)
 -\log\left(1+\frac1m\sum_{u=t-2m+1}^{t-m}N_{iju}\right).
\end{align*}
Never-observed activity needs its own indicator; its recency is not an invented zero. Count lookbacks and amount lookbacks should be named separately. Availability lags matter: a March statement finalized on April 10 was not available for an April 1 decision.

Industry and geography can enter as regularized categorical effects. Client names mainly support entity resolution; free-form names are not reliable economic predictors. Addresses can supply geography and parent--subsidiary matches, with confidence and provenance. Any target encoding must be fitted inside the training fold. An industry mean computed using future conversions leaks labels. Web announcements can add dated financing, expansion, hiring or acquisition indicators, but a current website cannot be retrospectively treated as its state two years ago. Backtesting requires archived publication and availability dates, and uncertain entity matches require explicit treatment.


\section{Product adoption, boosting, and probabilistic global models}
\subsection{Li, Sun, and Wilcox: sequential acquisition of banking products}

\citet{li2005cross}, Equations (1)--(5), model household purchases of several banking products using a latent demand-maturity measure. Let $O_j$ denote the position of product $j$ on a maturity scale. In the paper's notation, the household state is
\begin{equation}
 D_{Mi,t-1}=\sum_{j=1}^J O_jD_{ij,t-1}
 \left(\lambda_1\operatorname{ACCTNBR}_{ij,t-1}
      +\lambda_2\operatorname{BAL}_{ij,t-1}
      +\lambda_3\operatorname{Holding}_{ij,t-1}\right).
 \label{eq:li-maturity}
\end{equation}
The paper defines $D_{ij,t-1}$ as a dummy for purchasing product $j$ in the preceding month, not as an indicator of currently owning it. The other terms measure cumulative purchases, balances and elapsed holding time. Thus the source construction weights these accumulated attributes by recent purchase activity. Substituting current ownership would change the model. In particular, without any preceding-month purchase the displayed source maturity measure is zero; it is not automatically a persistent stock of financial sophistication.

The latent purchase utility is
\begin{equation}
 U_{ijt}=\beta_i|O_j-D_{Mi,t-1}|+\gamma_{1ij}\operatorname{COMPET}_{ij}
  +\gamma_{2ij}\operatorname{OVERSAT}_i
  +\gamma_{3ij}\operatorname{SWIT}_{it}+\varepsilon_{ijt},
 \label{eq:li-utility}
\end{equation}
with $Y_{ijt}=\mathbf1_{\{U_{ijt}>0\}}$ and contemporaneous errors $\varepsilon_{it}\sim N_J(0,\Sigma)$. The competition variable records opening the product category at another bank in the previous six months. Satisfaction comes from a survey. The source switching-cost indicator combines household characteristics; it is not a directly measured commercial switching expense. These details matter because this bank does not currently observe competitor holdings, and consumer demographic proxies should not simply be imported into a corporate model.

A negative $\beta_i$ makes purchase utility smaller when the product position is far from the household's maturity. Correlated errors allow contemporaneous demand for multiple products. The paper normalizes utility error scales and fixes one product position; these are necessary conventions for latent quantities, rather than directly observed dollar values. An empirical implementation must still examine whether the chosen data and restrictions distinguish the remaining parameters.

\subsubsection{From latent utility to a likelihood}

Write $\mu_{it}=(\mu_{i1t},\ldots,\mu_{iJt})$ for the deterministic utilities. For one product with unit error variance,
\[
 \Prb(Y_{ijt}=1)=\Prb(\varepsilon_{ijt}>-\mu_{ijt})
 =1-\Phi(-\mu_{ijt})=\Phi(\mu_{ijt}),
\]
where $\Phi(v)=\int_{-\infty}^v(2\pi)^{-1/2}e^{-z^2/2}dz$. For several products, let $\mathcal R(y)$ be the rectangle with coordinate $(0,\infty)$ when $y_j=1$ and $(-\infty,0]$ when $y_j=0$. The joint observation probability is
\begin{equation}
 p(y_{it}\mid\mu_{it},\Sigma)
 =\int_{\mathcal R(y_{it})}
 \frac{\exp\{-\tfrac12(u-\mu_{it})^{\mathsf T}\Sigma^{-1}(u-\mu_{it})\}}
 {(2\pi)^{J/2}|\Sigma|^{1/2}}du.
 \label{eq:li-probit}
\end{equation}
This explains why multiplying the $J$ univariate probabilities is generally wrong: that multiplication would impose a diagonal $\Sigma$. Conditional on household coefficients and the stated temporal model, multiply these rectangle probabilities over months and households. Integrate household coefficients against their population distributions to obtain the marginal likelihood.

The paper makes coefficients depend on household characteristics through a hierarchical model. In compact notation this has the form $\theta_i=Bz_i+e_i$, with normal individual deviations. Conditioning on $z_i$ shares information across similar households while retaining individual variation. For a corporate adaptation, $z_i$ would contain business attributes such as industry and size; this is a new specification requiring validation, not an asserted equivalence to household demographics.

\subsubsection{A complete estimation construction, with a source qualification}

The published article says that detailed priors and the simulation procedure are available on request. It does not contain enough detail to reproduce the authors' exact sampler from the article alone. The following construction therefore explains how to estimate the displayed model without that missing material; its priors and sampler are an explicit reconstruction, not a claim about the authors' unpublished choices.

One complete prior construction writes $\theta_i\mid B,D,R_\theta\sim N(Bz_i,DR_\theta D)$, where $D=\operatorname{diag}(\sigma_1,\ldots,\sigma_d)$. Give $\operatorname{vec}(B)$ and $\log\sigma_k$ proper normal priors with specified means and positive variances. Give free product locations and maturity coefficients proper normal priors, fixing the first product location as stated above. Give $R_\theta$ and the utility correlation matrix $\Sigma$ uniform densities over their positive-definite correlation-matrix domains; these priors are proper because those domains have finite volume in their free entries. Prior scales must be selected relative to standardized covariates and examined for sensitivity. These are reconstruction choices, not the authors' reported priors.

Writing $\psi$ for the population parameters, the resulting posterior is
\[
 p(\psi,\{\theta_i,u_{it}\}\mid y)\propto p(\psi)
 \prod_i\left[p(\theta_i\mid\psi)\prod_t
 \phi_J\{u_{it};\mu_{it}(\theta_i,\psi),\Sigma\}
 \mathbf1_{\{u_{it}\in\mathcal R(y_{it})\}}\right].
\]
Here $\phi_J$ is the $J$-variate normal density. This specification defines a posterior without supplying the unavailable original sampler or establishing identification from short histories.

Latent utilities permit a useful conditional update. Partition a multivariate normal into coordinate $j$ and the others. Completing its quadratic form shows
\begin{align*}
 U_j\mid U_{-j},\theta&\sim N(m_j,v_j)\text{ restricted by }Y_j,\\
 m_j&=\mu_j+\Sigma_{j,-j}\Sigma_{-j,-j}^{-1}(U_{-j}-\mu_{-j}),\\
 v_j&=\Sigma_{jj}-\Sigma_{j,-j}\Sigma_{-j,-j}^{-1}\Sigma_{-j,j}.
\end{align*}
For bounds $(l,b)$, draw $q$ uniformly between $\Phi((l-m_j)/\sqrt{v_j})$ and $\Phi((b-m_j)/\sqrt{v_j})$, then set $U_j=m_j+\sqrt{v_j}\Phi^{-1}(q)$. This is inverse-cdf sampling from the required truncated normal. Numerically extreme bounds require stable tail routines.

Conditional on product positions and maturity parameters, utilities are linear in the individual regression coefficients. Stack them as $u=X\theta+\epsilon$, with covariance $\Omega$, and suppose $\theta\sim N(m_0,V_0)$. Completing the square yields
\[
 V_1=(V_0^{-1}+X^{\mathsf T}\Omega^{-1}X)^{-1},\qquad
 m_1=V_1(V_0^{-1}m_0+X^{\mathsf T}\Omega^{-1}u),
\]
so $\theta\mid u\sim N(m_1,V_1)$. Hierarchical means admit analogous conditional Gaussian updates. Product positions and $\lambda$ enter through absolute values and products and generally lack conjugate updates. A proposal $\psi'$ from density $q(\psi'\mid\psi)$ can be accepted with probability
\[
 \min\left\{1,\frac{p(\psi'\mid u,\text{others})q(\psi\mid\psi')}
 {p(\psi\mid u,\text{others})q(\psi'\mid\psi)}\right\}.
\]
The conditional posterior in this ratio retains every factor involving the parameter being updated, including the individual-coefficient density when a population covariance changes. Its normalizing constant cancels. The same rule applies to a covariance parameterization constrained to remain positive definite, with transformation Jacobians when needed. Posterior purchase probabilities average Equation~\eqref{eq:li-probit} over retained draws. Convergence, prior sensitivity, identifiability and chronological prediction must be checked before treating this reconstruction as fitted evidence.

The source application observes eight products for 1,201 households over twelve months in 1997--1998 and evaluates prediction with held-out households. That provides a concrete consumer-banking application to short panels. It does not establish temporal robustness across macroeconomic regimes or effectiveness of contacting corporate clients. Its strongest contribution here is the hypothesis that past products and relationship development inform future product demand.

The initial proposal uses observed ownership, recent adoption, balances and interactions directly. A single latent product order is an optional challenger. Corporate clients can acquire products in different orders because of an acquisition, tender, import contract or financing maturity. A model that forces all clients onto one scale could miss those situations. No unavailable competition indicator will be silently coded as ``no competitor.''

\subsection{Friedman: gradient boosting from a loss function}

\citet{friedman2001greedy}, Algorithm 1, Section 4.5 and Section 5, constructs a prediction function as a sum of small fitted functions. Let $F(x)$ be the score and $L(y,F)$ its loss. Initialize with the best constant $F_0=\argmin_c\sum_i L(y_i,c)$. At iteration $m$, calculate the negative derivatives
\begin{equation}
 r_{im}=-\left.\frac{\partial L(y_i,F(x_i))}{\partial F(x_i)}\right|_{F=F_{m-1}}.
 \label{eq:pseudoresponse}
\end{equation}
Fit a base learner $h_m$ to these pseudo-responses, choose a step size by minimizing the original loss along $F_{m-1}+\rho h_m$, and update. Why regress on derivatives? A first-order expansion gives
\[
 \sum_i L(y_i,F_i+\epsilon h_i)
 =\sum_i L(y_i,F_i)-\epsilon\sum_i r_i h_i+O(\epsilon^2).
\]
Among unit-length changes at the observed points, alignment with $r$ gives the steepest local decrease. Fitting a small tree approximates this direction within the restricted class of tree functions. The procedure is greedy; it does not solve a global optimization over all possible tree ensembles.

For bank event indicators $y\in\{0,1\}$, use full log odds $\eta$ and Bernoulli loss
\[
 L(y,\eta)=\log(1+e^\eta)-y\eta,\quad
 p=\sigm(\eta),\quad L'=p-y,\quad L''=p(1-p).
\]
The initial score is $\log[\bar y/(1-\bar y)]$ when both classes occur. A regression tree partitions observations into leaves $R_{lm}$. In leaf $l$, the exact one-dimensional step minimizes
\[
 Q_l(\gamma)=\sum_{i\in R_{lm}}L(y_i,\eta_{i,m-1}+\gamma).
\]
Its derivative is $\sum_{i\in R_l}[\sigm(\eta_i+\gamma)-y_i]$. One Newton step from zero gives
\begin{equation}
 \gamma_{lm}\approx
 \frac{\sum_{i\in R_{lm}}(y_i-p_i)}{\sum_{i\in R_{lm}}p_i(1-p_i)},\qquad
 \eta_m(x)=\eta_{m-1}(x)+\nu\sum_l\gamma_{lm}\mathbf1_{\{x\in R_{lm}\}}.
 \label{eq:boosting-newton}
\end{equation}
Here $0<\nu\le1$ is shrinkage. This is a Newton approximation to the leaf minimizer, not an identity for the exact line search. Leaves containing only one class can have an unbounded unregularized optimum; minimum support, an explicit penalty or a bounded optimization rule must be specified and validated.

The paper's logistic derivation uses $y'\in\{-1,1\}$ and a half-log-odds score $F$, so $\eta=2F$. Its loss is $\log(1+e^{-2y'F})$ and its negative derivative is $2y'/(1+e^{2y'F})$. Transforming to $y=(y'+1)/2$ yields $2(y-p)$, while the second derivative becomes $4p(1-p)$. The Newton step in $F$ is therefore half the step in $\eta$. Keeping these codings distinct prevents a factor-of-two error when implementing the source algorithm.

Tree depth controls interaction complexity; number of trees and shrinkage jointly control fit. The paper's simulations and regression/classification examples illustrate these trade-offs, not a universal tuning choice. For the bank, boosting can discover interactions such as industry by currency exposure by recent activity. It is a challenger to regularized logistic hazards. It does not resolve missing competitor data, prove causality, or guarantee calibrated probabilities. Its additional complexity must earn a measurable improvement under chronological and segment-level evaluation.

\subsection{DeepAR: a shared recurrent probability model}

\citet{salinas2020deepar}, Section 3 and its training and prediction algorithms, model a conditional distribution for each related time series using shared recurrent-network parameters. The local copy is the 2019 arXiv manuscript; the bibliography records the 2020 journal publication. For series $i$, observations $z_{it}$, and covariates $x_{it}$, the chain rule factorization is
\begin{equation}
 p(z_{i,t_0:T}\mid z_{i,1:t_0-1},x_i)
 =\prod_{t=t_0}^{T}p(z_{it}\mid\theta_{it}),\quad
 h_{it}=f_\Theta(h_{i,t-1},z_{i,t-1},x_{it}),\quad
 \theta_{it}=g_\Theta(h_{it}).
 \label{eq:deepar}
\end{equation}
All clients share $\Theta$; different histories produce different states $h_{it}$. During training, the state receives observed previous values. During forecasting, future values are unknown and must be sampled recursively from the predictive distribution. This distinction is essential for multi-step evaluation.

\subsubsection{An explicit recurrent state}

For clarity, a standard long short-term memory state can be written without treating the network as an unexplained box. Let $v_t$ concatenate the previous scaled observation, covariates and previous hidden state. With trainable matrices and offsets,
\begin{align*}
 I_t&=\sigm(W_Iv_t+b_I),& F_t&=\sigm(W_Fv_t+b_F),\\
 O_t&=\sigm(W_Ov_t+b_O),& G_t&=\tanh(W_Gv_t+b_G),\\
 c_t&=F_t\odot c_{t-1}+I_t\odot G_t,&h_t&=O_t\odot\tanh(c_t).
\end{align*}
Products are coordinatewise. The forget gate $F_t$ controls persistence, $I_t$ controls new information and $O_t$ controls the exposed state. These equations specify one standard recurrent unit for the source model; their hidden dimension and regularization are empirical choices, not consequences of having a million clients.

For count data, output a mean $\mu_t>0$ and dispersion $\alpha_t>0$, with $k_t=1/\alpha_t$ in Equation~\eqref{eq:nb}. Positive outputs can be obtained using $\operatorname{softplus}(z)=\log(1+e^z)$, evaluated stably. The log-likelihood contribution is
\begin{align}
 \ell_t={}&\log\Gamma(z_t+k_t)-\log\Gamma(k_t)-\log\Gamma(z_t+1)\nonumber\\
 &+k_t\log\frac{k_t}{k_t+\mu_t}
  +z_t\log\frac{\mu_t}{k_t+\mu_t}.
 \label{eq:deepar-nb}
\end{align}
This gives $\E z_t=\mu_t$ and variance $\mu_t+\alpha_t\mu_t^2$. For real-valued data a Gaussian likelihood instead contributes $-\tfrac12\log(2\pi\sigma_t^2)-(z_t-\mu_t)^2/(2\sigma_t^2)$. A Gaussian permits negative outcomes, so it should not be applied uncritically to strictly positive balances or amounts.

The source rescales each series by a positive scale $\nu_i$ based on its conditioning history. Counts are fed to the recurrence as $z/\nu_i$, while the negative-binomial output mean is rescaled by $\nu_i$; its dispersion is scaled by $1/\sqrt{\nu_i}$ in the source parameterization. Scaling must use only the conditioning history for each training window. Using the future mean would leak the target. The source's sampling scheme also changes the emphasis across series; a bank must choose whether its objective weights clients, transactions, dollar margin or some explicit combination.

\subsubsection{Training and prediction, step by step}

Choose historical conditioning and prediction windows whose labels are fully observed. Run the recurrence through the conditioning window to establish its state, then minimize the sum of negative prediction-window log-likelihoods. For fixed dispersion, differentiating Equation~\eqref{eq:deepar-nb} gives
\[
 \frac{\partial\ell_t}{\partial\mu_t}
 =\frac{z_t}{\mu_t}-\frac{z_t+k_t}{\mu_t+k_t}.
\]
The chain rule propagates this derivative through the output map, softplus and recurrent states. More generally, for state $s_t=f_\Theta(s_{t-1},v_t)$ and total loss $\sum_t L_t$, the backward derivatives satisfy
\[
 a_t=\frac{\partial L_t}{\partial s_t}
      +\left(\frac{\partial s_{t+1}}{\partial s_t}\right)^{\mathsf T}a_{t+1},\qquad
 \nabla_\Theta L=\sum_t\left[
 \frac{\partial L_t}{\partial\Theta}\bigg|_{s_t}
 +\left(\frac{\partial s_t}{\partial\Theta}\bigg|_{s_{t-1}}\right)^{\mathsf T}a_t\right].
\]
The terminal derivative contains only the last loss. This is backpropagation through time: later errors inform parameters that created earlier states. Automatic differentiation evaluates these derivatives; it does not choose a valid train/test split or appropriate capacity.

At a forecast origin, compute the final observed state. For each simulation path, sample $z_{t+1}$ from its predicted distribution, feed that sample back to obtain the next state, and repeat through $t+H$. Path averages estimate future counts; path proportions with at least one qualifying event estimate event probabilities; empirical quantiles summarize uncertainty. Future covariates must be known, explicitly forecast or sampled under scenarios. Independently generated client paths do not capture unexplained bank-wide common shocks, so their sum does not automatically provide a calibrated aggregate-risk distribution.

The source evaluates probabilistic forecasting on several nonbank time-series datasets and gives architecture and training details in its supplement. Those results show a method for learning shared dynamics and predictive distributions. They do not establish that recurrent networks outperform well-tuned tabular models on 36 monthly observations. For this bank, DeepAR is justified as a bounded challenger when repeated counts and nonlinear dynamics matter. Its extra tuning, data weighting, seed variation, calibration and computational cost must be measured on this task before promotion.


\section{From prediction to profitable outreach}
\subsection{Define the intervention before estimating its value}

Forecasting who will transact and deciding whom to contact answer different questions. A relationship manager may already know which firms have imminent transactions. A model that predicts those transactions accurately can still waste sales capacity if the transactions would arrive without a call. Conversely, a client with modest natural propensity may be highly responsive to an appropriate offer.

Let $W=1$ mean assignment to a defined outreach protocol and $W=0$ mean assignment to a specified comparison, such as usual service without the additional product call. Let $M(w)$ be contribution margin over the outcome horizon under assignment $w$, before the separately recorded outreach cost. Define
\[
 m_w(x)=\E[M(w)\mid X=x],\qquad \tau(x)=m_1(x)-m_0(x),\qquad v(x)=\tau(x)-k(x).
\]
This is an intention-to-treat contrast: failed contact attempts and incomplete delivery are part of the outcome of assigning the protocol. Multiplying this $\tau$ by a contact-success probability would count delivery failure twice. If the desired effect is instead conditional on successful contact, successful delivery is a post-assignment selection event and needs a different identification argument.

The horizon, product, offer content, response window and usual-service condition must be explicit. The initial experiment should use one assignment per client group within the outcome horizon, with follow-up contacts governed by the assigned protocol. Repeated adaptive assignments require a longitudinal treatment model; the single-assignment formulas below do not automatically identify their value. A five-business-day call window can follow a monthly scoring cutoff, but the underlying transaction information may already be weeks old. A monthly model cannot promise detection within a few days of an unobserved transaction. Margin should include the finance-approved costs of funding, expected credit loss and implementation where relevant. Product notional and gross revenue are not substitutes for contribution margin. Sales labor cost is subtracted once, through $k$.

\subsection{Gutierrez and Gérardy: estimating heterogeneous response}

\citet{gutierrez2017causal}, Sections 2--3, review uplift methods for estimating differences in treatment response. For a binary outcome $Y$, the conditional average treatment effect is $\tau_Y(x)=\E[Y(1)-Y(0)\mid X=x]$. The corresponding derivations hold for a continuous margin outcome when the required moments exist.

The identification assumptions are consistency, conditional exchangeability and overlap. Consistency means the observed outcome is $Y=Y(W)$ for the defined intervention. Exchangeability means $(Y(0),Y(1))\perp W\mid X$. Overlap means $0<e(x)<1$, where $e(x)=\Prb(W=1\mid X=x)$. Under these conditions,
\begin{align}
 \E[Y(1)\mid X=x]
 &=\E[Y(1)\mid W=1,X=x]=\E[Y\mid W=1,X=x],\\
 \E[Y(0)\mid X=x]
 &=\E[Y(0)\mid W=0,X=x]=\E[Y\mid W=0,X=x].
 \label{eq:uplift-identification}
\end{align}
The first equality in each line uses exchangeability and the second consistency. Subtracting identifies the CATE. A record of past contacts establishes neither assumption. A banker's unrecorded knowledge of an impending deal can affect both treatment and outcome.

\subsubsection{Separate models, a joint model, and transformed outcomes}

The two-model approach fits $\widehat m_1$ on contacted observations and $\widehat m_0$ on controls, then uses their difference. Its strength is flexibility; its weakness is that two prediction errors enter the difference. If their conditional errors are $\epsilon_1$ and $\epsilon_0$, the effect error is $\epsilon_1-\epsilon_0$ and its variance is $\operatorname{Var}\epsilon_1+\operatorname{Var}\epsilon_0-2\operatorname{Cov}(\epsilon_1,\epsilon_0)$. Accurate marginal prediction does not guarantee accurate treatment-effect ranking.

A single joint model fits $m(x,w)$ and subtracts predictions at $w=1$ and $w=0$. For example, $m(x,w)=g(x)+w\theta(x)$ directly parameterizes heterogeneity for a continuous outcome. With a logistic binary-outcome model, the probability difference is $\sigm(g(x)+\theta(x))-\sigm(g(x))$, not the coefficient $\theta(x)$ itself. A model with no treatment interactions imposes strong restrictions on the pattern of uplift.

A transformed outcome has conditional mean equal to the effect:
\begin{equation}
 Y^*=\frac{W-e(X)}{e(X)(1-e(X))}Y
     =\frac{WY}{e(X)}-\frac{(1-W)Y}{1-e(X)}.
 \label{eq:transformed-outcome}
\end{equation}
Conditioning on $X=x$ and summing over treatment gives
\[
 \E[Y^*\mid x]=e\frac{m_1}{e}-(1-e)\frac{m_0}{1-e}=m_1-m_0.
\]
Thus least-squares regression of $Y^*$ on $X$ targets the CATE. But its second moment contains $\E[Y^2\mid W=1,x]/e+\E[Y^2\mid W=0,x]/(1-e)$, which becomes large near the propensity boundaries. Trimming changes the population being studied. Clipping denominators changes the estimator and generally introduces bias; neither is a free correction for absent overlap.

Uplift trees choose partitions for treatment heterogeneity instead of ordinary outcome prediction. A transparent split criterion can be derived from a squared-error fit to an effect signal. If two child nodes contain $n_L,n_R$ observations and effect means $\widehat\tau_L,\widehat\tau_R$, the reduction from fitting two means instead of one is
\[
 \frac{n_Ln_R}{n_L+n_R}(\widehat\tau_L-\widehat\tau_R)^2.
\]
To verify it, use the parent mean $(n_L\widehat\tau_L+n_R\widehat\tau_R)/(n_L+n_R)$ in the between-group sum of squares and expand. This illustrates the splitting principle surveyed in the paper; it is not a claim that every reviewed uplift-tree algorithm uses this exact criterion. Both treatment arms need support in each leaf, and independent estimation or validation is needed to avoid selecting spurious heterogeneity.

For binary outcomes one can imagine four individual types: responders only under contact, responders under either condition, nonresponders under either condition, and clients harmed by contact. The two marginal response probabilities do not identify these individual types because the joint distribution of $(Y(1),Y(0))$ is unobserved. The model estimates average response differences for comparable clients, not which counterfactual type a particular named firm certainly belongs to.

\subsection{Athey and Wager: doubly robust scores and policy choice}

\citet{wager2021policy}, their binary-treatment score and Section 2.3, estimate policy value using influence-function-based treatment scores and constrain policy complexity. Their full theory includes assumptions on sampling, nuisance estimation and policy classes. The following derivation states the binary-treatment calculation used here; it does not import an independent-sample asymptotic guarantee into a dependent bank panel.

With observed margin $M$ and fitted outcome means $\widehat m_w$, define
\begin{equation}
 \widehat\Gamma=\widehat m_1(X)-\widehat m_0(X)
 +\frac{W}{\widehat e(X)}[M-\widehat m_1(X)]
 -\frac{1-W}{1-\widehat e(X)}[M-\widehat m_0(X)].
 \label{eq:dr-score}
\end{equation}
The first two terms are a regression estimate of the difference. The weighted residuals correct systematic prediction errors using observed outcomes in each treatment arm. $\widehat\Gamma$ is a noisy effect score; it is not an observed individual causal effect and should not be used as a raw client ranking without estimating its conditional mean or a policy.

To establish double robustness, condition on $X=x$ and on nuisance functions estimated on separate data. Let $e$ and $m_w$ denote the true propensity and means. Expanding the conditional expectation gives
\begin{align*}
 \E[\widehat\Gamma\mid x]
 &=\widehat m_1-\widehat m_0
   +\frac{e}{\widehat e}(m_1-\widehat m_1)
   -\frac{1-e}{1-\widehat e}(m_0-\widehat m_0).
\end{align*}
Put $\delta_e=\widehat e-e$ and $\delta_w=\widehat m_w-m_w$. Subtracting $\tau=m_1-m_0$ and collecting terms yields
\begin{equation}
 \E[\widehat\Gamma\mid x]-\tau(x)
 =\delta_e\left[\frac{\delta_1}{\widehat e}
              +\frac{\delta_0}{1-\widehat e}\right].
 \label{eq:dr-bias}
\end{equation}
If the propensity is correct, this bias is zero even with misspecified outcome models. If both outcome means are correct, it is zero even with a misspecified propensity. Small errors enter as products when denominators are bounded away from zero. This does not protect against unmeasured confounding, outcome leakage, treatment misdefinition or lack of support: those failures invalidate the identification step before this algebra applies.

Cross-fitting separates nuisance fitting from score evaluation. Partition independent assignment units, fit nuisances outside each fold, and compute scores within the held-out fold. Related legal entities and repeated observations must stay together where appropriate. In a forecasting exercise, nuisance and policy training must also respect time; later outcomes cannot be used to build a model claimed to have existed earlier. Ordinary random cross-fitting of client-month rows is insufficient for either requirement. A prospective experiment with a known assignment probability removes the need to estimate that propensity, but still requires valid outcomes and policy evaluation.

\subsubsection{From treatment scores to a policy objective}

For a decision rule $\pi(x)\in\{0,1\}$, value relative to no additional contact is
\begin{equation}
 V(\pi)=\E[\pi(X)\{\tau(X)-k(X)\}],\qquad
 \widehat V(\pi)=\frac1n\sum_{i=1}^n\pi(X_i)(\widehat\Gamma_i-k_i).
 \label{eq:policy-learning}
\end{equation}
One can search a restricted class $\Pi$ for the rule maximizing this estimate, subject to budget and eligibility. The source sometimes expresses the objective with $(2\pi-1)\Gamma$. Since $\E[(2\pi-1)\Gamma]=2\E[\pi\Gamma]-\E\Gamma$, it has the same maximizer when the last term does not depend on $\pi$. Costs must be added consistently to whichever representation is used.

A short derivation explains why policy-class complexity matters. Let $\pi^*$ maximize true value in $\Pi$ and $\widehat\pi$ maximize estimated value. Add and subtract estimated values:
\begin{align*}
 V(\pi^*)-V(\widehat\pi)
 &=V(\pi^*)-\widehat V(\pi^*)
 +\widehat V(\pi^*)-\widehat V(\widehat\pi)
 +\widehat V(\widehat\pi)-V(\widehat\pi)\\
 &\le 2\sup_{\pi\in\Pi}|\widehat V(\pi)-V(\pi)|.
\end{align*}
The middle term is nonpositive by empirical maximization. In a simplified setting with independent observations, unbiased known scores, finite $\Pi$, and $|\Gamma-k|\le G$, Hoeffding's inequality and a union bound give
\[
 \sup_{\pi\in\Pi}|\widehat V(\pi)-V(\pi)|
 \le G\sqrt{\frac{2\log(2|\Pi|/\delta)}{n}}
\]
with probability at least $1-\delta$. This is an explanatory finite-class bound, not the paper's more general regret theorem. It shows how a policy can overfit noisy effects even when the scoring formula is unbiased. Nuisance error, dependent assignments and heavy-tailed margin require additional treatment.

For multiple mutually exclusive actions $a$, define $e_a(x)=\Prb(W=a\mid x)$ and $m_a(x)=\E[M\mid W=a,x]$. Relative to action $0$, the analogous score is
\[
 \Gamma_a=\widehat m_a-\widehat m_0
 +\frac{\mathbf1_{\{W=a\}}}{\widehat e_a}(M-\widehat m_a)
 -\frac{\mathbf1_{\{W=0\}}}{\widehat e_0}(M-\widehat m_0).
\]
Expanding by assignment as above establishes its corresponding robustness. Each action needs overlap. A separate binary analysis that treats other active offers as ``no contact'' changes the comparison and can be misleading.

\subsection{Elkan and Noto: learning from selectively observed positives}

\citet{elkan2008pu}, Lemma 1 and Sections 2--3, distinguish the true class $D\in\{0,1\}$ from the indicator $S$ that a positive is labeled. There are no labeled negatives: $\Prb(S=1\mid D=0,X)=0$. Under selected-completely-at-random labeling (SCAR),
\[
 \Prb(S=1\mid D=1,X=x)=c,\quad 0<c\le1.
\]
Writing $f(x)=\Prb(D=1\mid x)$ and $g(x)=\Prb(S=1\mid x)$, the law of total probability gives
\begin{equation}
 g(x)=c f(x),\qquad f(x)=g(x)/c.
 \label{eq:pu-lemma}
\end{equation}
This identity is correct under the assumptions. If $c$ is constant but unknown, $g$ preserves the ranking of $f$; it does not supply its calibrated level. If $c$ is known, the probability that an unlabeled example is truly positive is
\begin{align}
 w(x)&=\Prb(D=1\mid S=0,x)
 =\frac{\Prb(S=0\mid D=1,x)\Prb(D=1\mid x)}{\Prb(S=0\mid x)}\nonumber\\
 &=\frac{(1-c)f(x)}{1-cf(x)}
 =\frac{(1-c)g(x)}{c[1-g(x)]}.
 \label{eq:pu-weight}
\end{align}
For any classification loss $L(a,D)$, the conditional expected loss for an unlabeled observation is $wL(a,1)+(1-w)L(a,0)$. Labeled observations contribute $L(a,1)$. Averaging these terms over the observed $(X,S)$ distribution recovers the complete-data expected loss by iterated expectation, provided the probabilities used are correct. This explains the weighted-example construction without treating every unobserved bank transaction as a negative latent need.

\subsubsection{A limitation of the paper's proposed labeling-rate estimator}

On page 214 the paper proposes estimating $c$ by averaging $g(X)$ over held-out labeled positives and argues that this equals $c$ when $g$ is exact. Its argument substitutes $\Prb(D=1\mid X=x)=1$ for an observed labeled-positive example. That substitution requires more than knowing that this particular observation is positive: it requires its features to determine the positive class.

Under SCAR with overlapping probabilistic classes, the probability limit of that average is instead
\begin{equation}
 \E[g(X)\mid S=1]
 =c\E[f(X)\mid D=1]
 =c\frac{\E[f(X)^2]}{\E[f(X)]}.
 \label{eq:pu-c-limit}
\end{equation}
The second equality uses that labeling is random among positives. The third follows from $\E[f(X)\mathbf1_{\{D=1\}}]=\E[f(X)\E(D\mid X)]=\E[f(X)^2]$. Equality to $c$ requires $f(X)=1$ almost surely among positives, or another sufficient restriction; SCAR alone does not provide it.

For a concrete counterexample, let $X$ carry no information, $\Prb(D=1)=0.4$ and $c=0.5$. Then $g(x)=0.2$ for everyone, including labeled positives. The proposed average converges to $0.2$, not $0.5$, and dividing $g$ by that average produces $1$, not the true $0.4$. This counterexample concerns the stated estimator justification under general overlapping classes. It does not invalidate the identity $g=cf$ or the weighted-loss derivation when $c$ is known. The distinction is important for calibrated bank probabilities.

\subsubsection{Why the bank cannot assume random capture}

For banking, let $D$ denote a market-wide need and $S$ an own-bank observed positive. Capture depends on relationship strength, product, geography and competitor choices. The plausible relation is $g(x)=c(x)f(x)$, with unknown $c(x)$. Even ranking can reverse: client A can have need probability $0.40$ and capture $0.20$, giving an observed probability $0.08$, while client B has need $0.15$ and capture $0.80$, giving $0.12$. Ranking own-bank events then puts B first although A has greater total-market need.

An external sample with measured need and bank capture, or a defensible bound on capture, can add information. A randomized sales offer identifies its effect on observed outcomes under the experimental design; it does not by itself identify total market need or competitor wallet share. Web enrichment can improve prediction, but cannot remove this nonidentification merely by adding more covariates. The proposed product therefore labels its probabilities as bank-observed opportunities and treats total-market demand as a separate research target.


\section{The proposed bank sales product}
\subsection{The proposed model and its intended outputs}

The recommended first system has three connected components: a pooled product-event forecast, an outreach-effect model, and a capacity-constrained allocation rule. These components should share timestamped features but retain distinct targets. The initial product is a ranked queue of eligible client--product opportunities, with the information cutoff, expected timing and a specified next action. Once a prospective experiment supports incremental-value estimation, the queue can be ordered by expected net contribution under capacity constraints.

For product $j$, the primary forecast is a regularized pooled discrete-time logistic hazard. Industry and product effects are partially pooled; recency, frequency, activity changes, balances, holdings, maturity dates and prior contacts enter as available features. Product interactions reflect the cross-selling literature without imposing a universal product sequence. For repeatedly transacted products, a hurdle count or positive-amount model estimates intensity conditional on activity. The combined expectation is
\begin{equation}
 \E[V_{ijt}^{(H)}\mid C_{it}]
 =p_{ijt}^{(H)}\E[V_{ijt}^{(H)}\mid V_{ijt}^{(H)}>0,C_{it}],
 \label{eq:hurdle-value}
\end{equation}
provided the event is defined exactly as $V>0$. If a qualifying event has a threshold or includes a nonmonetary adoption, the count and amount targets must be adjusted; the identity cannot be attached to mismatched labels.

A global boosted hazard is the first nonlinear challenger. DeepAR is a later challenger for products where the full count distribution has business value. Original BG/NBD enters only if precise event timestamps are recovered; otherwise the discrete analogue or simple recency/frequency features carry its useful insight. A latent maturity or multivariate-probit model is optional because its stronger product-order and data requirements are not established for corporate clients.

The forecast answers what happens under the contact and market behavior represented in its training data. If the bank materially changes its sales policy, that outcome distribution may change. A historical-event probability should not be mislabeled as an untreated probability. The causal layer models the intervention explicitly and is evaluated on the policy that the bank intends to use.

\subsection{Data construction and model fitting}

Resolve legal entities and relationship groups before generating labels. Keep both the legal-client identifier and a group identifier so that related subsidiaries can be coordinated, and so that a parent's information does not leak into a supposedly independent subsidiary holdout. Product eligibility is evaluated at the cutoff. Existing ownership removes a pair from first-adoption risk but can create a separate expansion opportunity. Missing feeds, unobservable periods and delayed settlement records have explicit states.

A practical training sequence is to create compact monthly aggregates, construct lag features using only past available data, build risk sets, and then fit one shared model per product or a multi-product model with product interactions. Thirty-six million client-month rows before eligibility and horizons is a large table, but not automatically a deep-learning problem. Partitioning by cutoff and product, reusing common features and streaming batches can keep computation manageable. Runtime and memory must be measured on a representative extract; no hardware requirement follows from row count alone.

If non-event rows are sampled for computational reasons, retain sampling probabilities. An inverse-probability-weighted loss targets the original population loss: for a sampled indicator $R$ with inclusion probability $s(x,y)$, $\E[R L/s\mid x,y]=L$. With outcome-only case/control sampling, the sample log odds satisfy
\[
 \logit p_{\rm sample}(x)=\logit p_{\rm population}(x)+\log(s_1/s_0).
\]
Thus subtracting $\log(s_1/s_0)$ corrects the intercept under this specific design. Arbitrary client-dependent sampling requires appropriate weighting or a fuller correction. Unweighted oversampling followed by interpreting scores as population probabilities is wrong.

Fit a calibration model on a later, disjoint calibration window. A simple logistic calibration is $p_{\rm cal}=\sigm(a+b\logit p)$, with $(a,b)$ estimated by Bernoulli likelihood. Perfect calibration would give $a=0,b=1$. An isotonic alternative fits a nondecreasing step function by minimizing squared probability error subject to ordered fitted values; adjacent blocks that violate monotonicity are pooled to their weighted mean until all blocks are ordered. Either method needs enough positive outcomes per product and independent final evaluation. Calibration cannot repair missing predictors or establish causal validity.

\subsection{Learning the value of an outreach assignment}

Use historical contact records first to assess support and data quality. Reconstruct who was eligible, who was assigned a call, who was actually reached, what offer was delivered, and when the outcome became observable. Estimate propensity and outcome models using cross-fitting at the appropriate entity-group level, with chronological separation for prospective claims. Regress the doubly robust scores on pre-assignment features, or fit a restricted policy directly to those scores. Individual score noise is not individual effect evidence.

Historical models remain provisional where banker's private information plausibly drives selection. A prospectively randomized pilot should cover the eligible population on which the future policy will operate, stratified by product, segment and propensity band. Randomize the additional outreach protocol against the agreed usual-service condition and record assignment probabilities. Analyze assignment as the treatment. Contact completion is a useful operational diagnostic; conditioning the primary causal analysis on completed contact would undo randomization.

Where overlap is absent, the system may still provide an event forecast. It should not extrapolate an experimental effect to unsupported actions as though it were measured. Product-specific eligibility restrictions can be respected in the experiment without pretending that excluded clients have comparable treatment effects. Sensitivity analyses can show how plausible unrecorded selection would change observational conclusions; they cannot prove exchangeability.

\subsection{Allocate expected contribution under actual constraints}

For action $a$ on client $i$, let $\widehat\tau_{ia}$ estimate incremental margin from assigning the complete protocol and $k_{ia}$ its expected outreach cost. Use
\begin{equation}
 \widehat v_{ia}=\widehat\tau_{ia}-k_{ia}.
 \label{eq:action-value}
\end{equation}
No extra transaction-probability or delivery-probability multiplier belongs in this expression because both are already integrated into the assignment effect. A different model that separately predicts response probability and conditional margin must derive its value from those components and use a consistent baseline.

At a monthly decision point, choose binary assignments $z_{ia}$ by
\begin{align}
 \max_z\quad &\sum_{i,a}\widehat v_{ia}z_{ia},\nonumber\\
 \text{subject to}\quad &\sum_{i,a}h_{ia}z_{ia}\le B,\nonumber\\
 &\sum_a z_{ia}\le1\quad\text{for each client }i,\nonumber\\
 &\sum_{i,a\text{ assigned to }r}h_{ia}z_{ia}\le B_r
       \quad\text{for each salesperson }r,\nonumber\\
 &z_{ia}=0\text{ for ineligible actions},\qquad z_{ia}\in\{0,1\}.
 \label{eq:mip}
\end{align}
Add territory, product-specialist capacity, contact-frequency and relationship-group coordination constraints as required. Costs in the objective and hours in the constraint serve different purposes: one values resources economically and the other prevents assigning more work than can be delivered. Bundle actions need their own incremental values when product effects interact; adding two separate effects assumes away those interactions.

With identical hours and no other constraints, selecting the largest positive values up to capacity is exact. With unequal hours, value-per-hour sorting is only a heuristic for the indivisible knapsack problem. For example, with capacity four hours, an action worth five units requiring three hours has higher value per hour than each of two actions worth three units requiring two hours, yet selecting the two smaller actions yields six instead of five. An integer optimizer or a documented approximation should handle the actual constraints. The optimizer's solution is only as credible as its estimated values and constraints; it is not evidence of global economic optimality.

\subsubsection{An illustrative queue}

Table~\ref{tab:queue} uses hypothetical numbers to show why event propensity and outreach value should both be visible. Margin effects already include the possibility of unsuccessful contact. With three available hours, B and C yield estimated net contribution $650$, exceeding C and D at $600$. Client A has the highest event probability but a negative estimated net outreach value.

\begin{table}[htbp]
\centering
\caption{Hypothetical assignments, in dollars and salesperson hours. No bank data were used.}
\label{tab:queue}
\begin{tabular}{lrrrrr}
\toprule
Client & Event probability & Incremental margin & Contact cost & Net value & Hours\\
\midrule
A & 0.80 & 40 & 50 & $-10$ & 1\\
B & 0.30 & 500 & 100 & 400 & 2\\
C & 0.25 & 300 & 50 & 250 & 1\\
D & 0.40 & 430 & 80 & 350 & 2\\
\bottomrule
\end{tabular}
\end{table}

Estimated effects have uncertainty. A conservative value such as $\widehat\tau-\lambda\widehat s-k$ can express a deliberate aversion to estimation risk, but $\lambda$ is a business preference and this score is not automatically a valid lower confidence bound for a selected portfolio. Selection among a million noisy estimates can amplify optimism. Independent policy evaluation, limits on unsupported extrapolation, and a reserved randomized exploration allocation are more informative than displaying narrow individual intervals without adequate data.

\subsection{Unusual activity as a secondary signal}

A separate anomaly score can highlight an observed change that may warrant a conversation. At cutoff $t$, compare the newly observed count $n_t$ with a predictive distribution fitted using information available before that observation. A right-tail score is $u_t=\Prb(N_t\ge n_t\mid C_{t-1})=1-F_{t\mid t-1}(n_t-1)$. For unusual declines, use the lower tail. A very small tail probability means the observation was surprising under that model; it does not establish financial distress, misconduct, product need or responsiveness to contact.

An anomaly can enter the next forecast as a lagged feature or trigger review together with an appropriate product signal. With a million clients, many tail observations arise even under a well-calibrated model. Thresholds should reflect the review capacity, expected false alerts and outcome validation. Do not mix this anomaly queue with a sales-value queue without measuring whether the additional calls help.

\subsection{The banker interface and the feedback needed for learning}

A recommendation should show the client and relationship group, product, eligible action, horizon, information cutoff, probability band, relevant magnitude forecast, owner and contact deadline. After causal validation it can add estimated incremental contribution and uncertainty. Three dated descriptive reasons might be recent cross-border payment growth, a known facility maturity and a compatible existing product. These are associations supporting the forecast, not claims that the feature causes a purchase.

The banker records acceptance, override reason, attempt, completion, offer and outcome. Recommendations, overrides and withheld randomized assignments remain in the history so that later evaluation can distinguish the model's decision from delivery behavior. Outcomes must also be collected for uncontacted clients; a dataset containing only contacted successes cannot support the intended comparison. Monthly scoring and daily CRM execution can coexist, provided the interface states the actual data age.

\subsection{Why this model addresses the stated problem}

Pooling learns common patterns from many clients while shrinkage limits the parameters inferred from each short history. Discrete-time hazards match the monthly observation frequency and distinguish first adoption from recurring events. Product histories use cross-selling information that the bank actually has. Count and amount models add economically relevant scale. The explicit bank-observed target respects the unobserved competitor relationships. The treatment model asks whether sales attention changes the outcome, and the allocation model respects the finite sales force.

These properties make the design capable of answering the requested decisions under stated assumptions. They are not a proof that it will improve sales. That conclusion requires the historical and prospective comparisons below, including simple strategies that a banker could implement without a complex model.


\section{Evaluation, operations, and governance}
\subsection{Chronological evaluation with fully matured labels}

A random split of client-month rows would allow future conditions, repeated client histories and overlapping outcome windows to leak across training and testing. The main test should instead reconstruct what would have been known at each historical decision date. Client-group holdouts answer the additional cold-start question, but they do not replace temporal holdouts for existing clients.

A concrete three-month-horizon design with 36 months is shown in Table~\ref{tab:time}. It is an example to adapt to event frequency and feature availability. Features at cutoff $t$ use only information available then; labels cover $t+1,t+2,t+3$. Hyperparameters and calibration are frozen after month 27. The final test uses cutoffs 27--33, whose outcomes mature by month 36. If the model is refitted during the test, the refitting rule must be fixed in advance and may use only labels that have matured at that historical date.

\begin{table}[htbp]
\centering
\caption{Illustrative chronology for a three-month horizon.}
\label{tab:time}
\begin{tabular}{P{0.18\linewidth}P{0.24\linewidth}P{0.45\linewidth}}
\toprule
Use & Forecast cutoffs & Information and labels\\
\midrule
Model fitting & Earliest feasible cutoff through 21 & Lookbacks end at each cutoff; all three-month labels end by 24.\\
Validation & 22--24 & Labels end by 27; select model complexity and fit calibration using this later information.\\
Final testing & 27--33 & Freeze choices at 27; forecast future outcomes through 36. No test label informs selection.\\
\bottomrule
\end{tabular}
\end{table}

The validation window is short and selecting many models on it can overfit. Within the earlier data, use a small number of rolling-origin folds to narrow choices, and reserve the final test for one comparison. One-month targets have different label maturity and should be evaluated separately. Three-month targets at successive cutoffs overlap, so even distinct test rows can share the same transaction. Uncertainty must account for that overlap. Seven test origins do not establish robustness across many economic cycles.

\subsection{Constructed practitioner baselines and model comparisons}

Before inspecting final labels, define simple strategies a relationship manager could reasonably use. For this problem the comparison set should include:
\begin{enumerate}[leftmargin=*]
 \item Current banker prioritization, reconstructed from its actual rules and available information, to establish the operational comparator.
 \item Random selection among eligible clients within product and territory, repeated using recorded seeds, to measure the value of targeting at the same capacity.
 \item A fixed recency/frequency rule: prioritize clients with more active months in the last three months, breaking ties by more recent activity. This directly tests whether a complex model adds to persistence.
 \item A product--industry base-rate rule estimated only on the fitting period, with sparse industries shrunk toward the product mean. This tests whether categories alone explain the ranking.
 \item A seasonal-repeat rule prioritizing clients with activity in the corresponding month of the previous year, with the product base rate for unavailable history. This tests a simple annual pattern while recognizing that only three years are observed.
 \item A known-event rule prioritizing documented facility maturities or product-renewal dates inside the horizon, using a fixed fallback for clients without such dates. This tests commercially obvious timing information.
\end{enumerate}
These rules are fixed before final testing. They are falsification comparisons, not a collection of final-test scores against which to tune the proposed model. Current-process records may themselves be incomplete; that limitation must be reported rather than replaced silently by a weak artificial baseline.

Compare a tuned regularized logistic model, a partially pooled hazard, a boosted hazard and, if warranted, DeepAR. Each must receive the same eligible population, historical information boundary and outcome definitions. Feature-rich nonlinear models should also be compared with feature-rich classical models so that extra data are not mistaken for an architectural benefit. Count models use the same count definition and forecast horizon; event models are compared on event probabilities.

Report results separately for new and established clients, sparse and frequent transactors, small and large exposures, industries, products, relationship groups and calendar periods. The salient situations should be chosen before final evaluation. A complex model losing to a simple rule in a situation it is intended to handle is a main finding and prevents promotion until understood or repaired. A small, uncertain difference is not evidence of superiority. Rare subgroups may require more data before any ranking can be defended.

\subsection{Forecast scores, calibration, and usable capacity}

For binary outcomes, average log loss is
\[
 -\frac1n\sum_i[y_i\log p_i+(1-y_i)\log(1-p_i)],
\]
and the Brier score is $n^{-1}\sum_i(y_i-p_i)^2$. To see why calibrated probabilities matter, suppose the true probability at a given feature value is $q$. Expected Brier loss is
\[
 q(1-p)^2+(1-q)p^2=q(1-q)+(p-q)^2,
\]
which is minimized at $p=q$. Expected log loss has derivative $-q/p+(1-q)/(1-p)$, also zero at $p=q$, with positive curvature. Thus both reward probability accuracy, not just ordering. Scores should be computed on the actual eligible population or appropriately weighted back to it.

For a capacity of $K$ recommendations, define precision at $K$ as $K^{-1}\sum_{i\in\operatorname{Top}K}y_i$. Recall at $K$ divides the same numerator by all positive outcomes. Lift at $K$ divides precision at $K$ by population event prevalence. Precision--recall curves are useful for rare events, while area under a receiver-operating curve can conceal poor precision at the small fraction the bank can contact. Report the operating point matching actual salesperson hours, as well as fixed-client-count comparisons when useful.

A calibration plot compares mean predicted probability with observed frequency in prespecified bins, with uncertainty and sample counts. Examine it by product, horizon and material segment. For count or amount forecasts, use predictive log scores and interval coverage together with errors in the economically relevant aggregate. Coverage without interval width can reward an uninformative distribution, and a good average error can hide failure on large exposures.

Historical transaction lift is not incremental sales lift. A forecast-only queue can be assessed for probability accuracy and concentration of future transactions. The primary commercial comparison for an outreach policy is incremental contribution at matched capacity. Forecast accuracy is an intermediate diagnostic; it cannot establish that calling the selected clients creates value.

\subsection{Uncertainty and policy evaluation}

Compare candidate and baseline on the same held-out observations, using paired differences. Resampling entire client groups preserves within-group histories. Resampling calendar blocks can address some common shocks and overlapping horizons. With only 36 months, neither procedure creates additional regimes; intervals for performance under a new macroeconomic environment remain weakly identified. Report both segment variation and calendar variation, and state when the number of blocks is too small for reliable asymptotic inference.

For a policy fixed before evaluation on independent units, Equation~\eqref{eq:policy-learning} estimates its incremental value. If $G_i=\pi(X_i)(\widehat\Gamma_i-k_i)$ and the assumptions support independent units, the usual standard error is $s_G/\sqrt n$. With clustered assignment, first aggregate to the randomization unit and use its sampling design. When comparing two fixed policies, use $(\pi_1(X_i)-\pi_0(X_i))(\widehat\Gamma_i-k_i)$ on the same observations, rather than treating their value estimates as independent. Reusing the evaluation data to choose the policy makes these simple intervals optimistic.

A portfolio selected at a fixed monthly capacity couples the client decisions. Evaluate the allocation procedure at its intended operating capacity, including ownership and territory restrictions. A population rule with unrestricted contact fraction and a monthly optimizer with a fixed staff budget are not automatically the same policy. Historical replay is informative only to the extent that the recorded actions and outcomes support its counterfactual comparisons.

\subsection{Two prospective experiments answer two different questions}

The first experiment estimates the effect of the outreach protocol. Randomize eligible client groups within declared strata to additional contact or usual service. Record outcomes for everyone, including unsuccessful attempts. Compare assignment groups over the complete outcome horizon. This establishes whether the action has value and supplies data for heterogeneous-effect estimation.

The second experiment evaluates allocation at equal resources. Assign comparable portfolios, territories or appropriately separated periods to the proposed queue or the current prioritization process, with equal salesperson hours and consistent offers. Choose the unit to limit contamination across relationship groups and bankers. Evaluate total contribution across each assigned portfolio, including clients not contacted, so that differences in who was selected do not become a post-treatment sample restriction. A switchback design needs enough periods and control of carryover; it is not automatically suitable for long-lived financing decisions.

The two experiments may be combined in a carefully designed trial, but their estimands remain distinct. Comparing contacted model-selected clients with uncontacted randomly chosen clients confounds selection and treatment. Comparing conversion rates only among successfully reached clients additionally conditions on delivery. Neither establishes the value of the allocation policy.

\subsubsection{Power and exposure must be planned from event rates}

For an illustrative equal-size randomized binary-outcome trial, let $p$ approximate the response probability in each arm, $\Delta$ the minimum detectable difference, and $n_a$ the number of independent units per arm. The variance of the difference is approximately $2p(1-p)/n_a$. Requiring a standardized signal of $z_{1-\alpha/2}+z_{1-\beta}$ gives
\begin{equation}
 n_a\approx\frac{2p(1-p)(z_{1-\alpha/2}+z_{1-\beta})^2}{\Delta^2}.
 \label{eq:power}
\end{equation}
Here $\alpha$ is the test's false-positive probability and $1-\beta$ its desired power. This is a normal-approximation planning formula, not a final rare-event design. For a continuous margin outcome with approximately common variance $\sigma_M^2$, replace $p(1-p)$ by $\sigma_M^2$ and express $\Delta$ in dollars. Heavy tails, unequal variances and product mix may require simulation-based planning.

If a cluster contains $m$ clients with common variance $\sigma^2$ and pairwise correlation $\rho$, the variance of its sum is $m\sigma^2+m(m-1)\rho\sigma^2$, giving the familiar variance multiplier $1+(m-1)\rho$. Unequal cluster sizes and stratification need their actual design calculation. One million customers therefore does not guarantee a well-powered trial: rare product events, limited outreach, few independent bankers and long outcome windows can be the binding constraints. Historical rates, costs and margin variability are required before setting the final sample size.

\subsection{Promotion criteria and operational failure conditions}

The forecast release should meet predeclared probability-calibration tolerances, outperform or remain acceptably competitive with the simple baselines at operating capacity, and show no unexplained failure in material segments. The causal release should demonstrate positive incremental value against the agreed comparison with uncertainty appropriate to the trial, together with acceptable contact burden and delivery performance. A model that passes a technical screen remains a candidate; it is not thereby proven superior.

Missing outcome maturity, leakage, invalid entity matching, absent treatment support or a broken randomization record invalidate the relevant comparison and require repair. Weak predictive performance rejects that candidate but need not reject the entire research direction. For example, a recurrent network failing to improve on recency does not invalidate the pooled logistic proposal. All reported comparisons should distinguish such failures from ordinary statistical uncertainty.

Monitor feature availability, score distributions, event prevalence, calibration, contact volume, delivery rates, incremental margin and segment coverage. Retraining should follow a documented schedule or drift investigation, using matured labels. Preserve the model version, data cutoff and recommended action for every decision. Feedback from banker overrides can improve the product, but an override is not automatically a correct label; it may reflect capacity, private information or preferences that require separate analysis.


\section{Limits, extensions, and implementation sequence}
\subsection{A staged implementation with concrete deliverables}

Begin with a small number of products whose outcomes and economics are well measured, for example one frequently transacted product and one first-adoption product. This tests both recurrence and adoption without trying to resolve every commercial definition simultaneously. The choice is a proposed scope assumption, to be made with product owners using data availability and potential value.

The first stage produces a signed-off outcome dictionary, timestamp and availability rules, entity-group mapping, eligibility logic, missing-data audit and a reproducible monthly panel. Reconcile sampled labels with source systems and CRM records. Establish whether exact transaction dates can be recovered, which determines whether original BG/NBD is a valid benchmark. Record known facility maturities separately from model-inferred signals.

The second stage produces the constructed simple-rule comparisons, a regularized logistic forecast and a partially pooled hazard, followed by a bounded boosting comparison. Its deliverables are chronological predictions, calibration estimates, segment results, paired uncertainty, runtime and memory measurements, and a proposed queue at actual capacity. Web enrichment should enter as an ablation: compare the same model and dates with and without the additional fields. A current web crawl cannot serve as historical data unless archived timestamps establish availability.

The third stage produces a prospectively randomized contact protocol, sample-size analysis, assignment log and complete outcome collection. Historical uplift estimates can nominate strata and hypotheses, but cannot replace the trial. After outcomes mature, estimate assignment effects and a restricted allocation policy using separate fitting and evaluation data. If the action itself has no demonstrated value, changing the ranking algorithm is not a sufficient remedy; reconsider the offer, timing or eligible population.

The fourth stage evaluates the proposed allocation against current prioritization at equal capacity and measures its adoption in the CRM workflow. Expand products, horizons or data sources only after the existing comparison supports the additional complexity. DeepAR and latent maturity models enter as optional research candidates when simpler methods leave a specific, measurable deficiency. The product can remain useful with a classical pooled model if that model meets the commercial criteria.

\subsection{A bounded research protocol before any model campaign}

This document does not execute model training, customer scraping or a sales experiment. Before those activities, the team should set the dataset version, financial outcome, baseline, primary criterion, conditional evaluation situations, uncertainty method, compute and attempt budget, and conditions for stopping or repairing a run. A practical comparison uses a ladder: simple rules, a tuned classical model, a plain nonlinear candidate, then any enhanced candidate. The final test remains untouched throughout tuning.

Every material choice is a hypothesis. Lookback lengths should reflect plausible business cycles and be tested on earlier chronological folds. Regularization controls unstable sparse effects and requires sensitivity checks. Tree depth and shrinkage trade interactions against overfit. A neural candidate needs its own capacity and optimizer search, multiple recorded seeds, early stopping based on an earlier validation window and a sufficient budget for meaningful downstream comparison. A convenient architecture from a retail paper is not a bank-specific default. If the available budget cannot support a fair neural comparison, defer that comparison and retain the validated simpler candidate.

Before expanding computation, run a focused check that the labels, feature timing, loss and probability units are correct. A small successful training run establishes execution, not scientific or commercial adequacy. Store exact configurations, data versions, seeds, runtime, CPU/GPU mode, complete predictions and failure records. A localized infrastructure failure warrants repair within the agreed scope; an invalid target or corrupted outcome record requires revising the experiment before continuing.

\subsection{What remains uncertain}

Three years can support cross-sectional learning but gives little evidence about rare economic regimes. Corporate product demand may be driven by acquisitions, tenders or financing plans absent from transaction histories. External announcements may improve those signals, but their timeliness, coverage and entity matching must be measured. Larger feature sets do not resolve an unobserved demand denominator.

Historical contacts may be too selectively recorded to identify effects. Randomization can address assignment confounding within the experimental population, while interference, incomplete follow-up and changing offers still require attention. Value measured over a short horizon can miss later attrition or credit losses. Where long-run margin cannot yet be observed, report the chosen short-run outcome and the unresolved long-run economics explicitly.

The recommendation is therefore to build the measurable bank-opportunity forecast first, collect credible intervention evidence alongside it, and allocate sales capacity using validated incremental value when that evidence becomes available. The mathematics specifies how to do this coherently. The data and prospective comparisons determine whether it succeeds for the bank.


\clearpage
\appendix
\section{Notation and derivation index}
\subsection{Notation and units}

\begin{longtable}{P{0.20\linewidth}P{0.73\linewidth}}
\toprule
Symbol & Meaning\\
\midrule
\endhead
$i,j,t$ & Business client, product, calendar-month cutoff.\\
$k,s,H$ & Elapsed month within a spell, spell index, forecast horizon in months.\\
$C_{it},X$ & Information available at the decision cutoff; $X$ is generic notation in method derivations.\\
$Y_{ijt}^{(H)}$ & Indicator of a defined own-bank event during the next $H$ months.\\
$N,V,M$ & Transaction count, monetary amount, contribution margin. These are distinct targets with different units.\\
$h,p$ & Monthly conditional event hazard and horizon event probability.\\
$W,M(w)$ & Assigned outreach action and margin under that assignment.\\
$e(x),m_w(x)$ & Treatment-assignment probability and conditional outcome mean for action $w$.\\
$\tau,v,k$ & Incremental margin before outreach cost, net incremental value, outreach cost. The symbol $k$ in a cost expression differs from elapsed time by context.\\
$z_{ia},h_{ia},B$ & Action assignment, required hours and available salesperson hours. Here $h_{ia}$ denotes hours, not the hazard $h_{ijt}$.\\
$D,S,c(x)$ & Latent market need, observed positive/capture indicator, probability of observing a positive need.\\
$\lambda,p,r,\alpha,a,b$ & BG/NBD active arrival rate, post-transaction dropout probability and gamma/beta population parameters. Its $p$ differs from the horizon-event probability.\\
$\mu,k,\alpha$ & Negative-binomial mean, shape and dispersion, with $\alpha=1/k$. These are local distribution parameters.\\
\bottomrule
\end{longtable}

\subsection{Where each paper enters the proposed product}

\begin{longtable}{P{0.23\linewidth}P{0.33\linewidth}P{0.35\linewidth}}
\toprule
Source & Reconstructed contribution & Bank use and boundary\\
\midrule
\endhead
\citet{montero2021global} & Finite-history representation and shared-class generalization calculation. & Pool clients; dependence and only 36 calendar periods prevent a naive sample-size guarantee.\\
\citet{singer1993time} & Hazard, survival, censoring likelihood and person-period logistic estimation. & Predict first adoption in monthly intervals with correct eligibility and risk sets.\\
\citet{willett1995multiple} & Multiple-spell likelihood and history-dependent risk clocks. & Handle repeated FX or transaction opportunities without treating adoption as a recurring first event.\\
\citet{fader2005counting} & Gamma/beta mixture likelihood, active probability and future-count expectation. & Recency/frequency benchmark if exact dates exist; a separately defined monthly analogue otherwise.\\
\citet{li2005cross} & Maturity-dependent multivariate probit, hierarchical heterogeneity and latent-utility estimation. & Motivate product-history features; competitor data and consumer order assumptions do not automatically transfer.\\
\citet{friedman2001greedy} & Functional gradient updates, logistic leaf steps and shrinkage. & Nonlinear hazard challenger, with probability calibration and chronological selection.\\
\citet{salinas2020deepar} & Shared recurrent likelihood, count outputs and recursive simulation. & Probabilistic count challenger; no presumption of benefit on short monthly bank histories.\\
\citet{gutierrez2017causal} & Treatment identification, separate/joint models, transformed outcomes and splitting principle. & Estimate average outreach response rather than ordinary conversion propensity.\\
\citet{wager2021policy} & Doubly robust score, bias cancellation and policy-value objective. & Learn and evaluate contact allocation under overlap, valid assignment assumptions and capacity.\\
\citet{elkan2008pu} & SCAR identity and weighted unlabeled loss; labeling-rate estimator limitation derived here. & Explain partial observation and why unknown capture does not identify total-market demand.\\
\bottomrule
\end{longtable}

\subsection{Source versions and derivation scope}

The bibliography gives complete publication metadata. The accompanying source manifest maps every citation to a downloaded full text, its retrieval address, version, checksum and technical reading anchors. Author manuscripts are retained when they are the accessible full text; in particular, the DeepAR manuscript has three named authors while the journal publication lists four. The downloaded copy is not represented as the final publisher layout. Scanned source papers remain unchanged, with extracted or OCR text retained separately for research use.

The main text supplies the likelihoods, losses, prediction formulas and estimation arguments used in the proposal. It does not reproduce unrelated empirical tables or claim to prove every theorem in each cited paper. The policy-learning finite-class bound is an explanatory derivation, not a reproduction of the source's general asymptotic theorem. The cross-selling estimation construction is explicitly supplied because the article does not publish all its sampler details. The monthly inactivity model, bank allocation example, capture counterexamples and experimental design are adaptations developed for this problem.

All numerical queue values are hypothetical. No bank records were supplied, no prediction model was trained and no commercial performance was measured. The proposed comparisons specify what evidence would be needed to move from a mathematically coherent design to a successful sales product.


\clearpage
\begingroup
\raggedright
\interlinepenalty=10000
\bibliographystyle{plainnat}
\bibliography{references}
\endgroup
\end{document}
